% =====================================================================
% Artigo — P2 Finanças do Setor Público — Matrícula 231019943
% Compilar: latexmk -pdf artigo_P2_231019943.tex  (ou pdflatex + bibtex + pdflatex x2)
% =====================================================================
\documentclass[article,11pt,oneside,a4paper,english,brazil,sumario=tradicional]{abntex2}
\usepackage[T1]{fontenc}
\usepackage[utf8]{inputenc}
\usepackage{lmodern}
\usepackage{microtype}
\usepackage{graphicx}
\usepackage{booktabs,makecell,threeparttable,float,longtable}
\usepackage{amsmath,amssymb}
\usepackage{icomma}
\usepackage{xcolor}
\usepackage[alf,abnt-emphasize=bf,abnt-etal-cite=3,abnt-etal-list=0]{abntex2cite}
\usepackage{xurl}
\usepackage{caption}
\captionsetup{font=small,labelfont=bf,justification=raggedright,singlelinecheck=false}
\setlrmarginsandblock{2.5cm}{2cm}{*}
\setulmarginsandblock{2.5cm}{2cm}{*}
\checkandfixthelayout
\setlength{\parindent}{1.25cm}
\setlength{\parskip}{0.15cm}
\SingleSpacing
\renewcommand{\baselinestretch}{1.15}
\hypersetup{colorlinks=true,linkcolor=blue!50!black,citecolor=blue!50!black,urlcolor=blue!50!black,
  pdftitle={Transferências intergovernamentais, ilusão fiscal e efeito flypaper: Manaus, Parintins e Espanha (2016–2025)},
  pdfauthor={João Pedro Furtado Ximenes da Rocha}}
% Gerado automaticamente por 06_codigos_e_planilhas/04_numeros_artigo.py — não editar à mão
\newcommand{\anoIniMan}{2016}
\newcommand{\depManMedia}{58,5}
\newcommand{\depManIni}{52,7}
\newcommand{\depManFim}{57,6}
\newcommand{\depManMin}{52,7}
\newcommand{\depManMax}{63,0}
\newcommand{\autManMedia}{25,9}
\newcommand{\autManIni}{24,4}
\newcommand{\autManFim}{28,0}
\newcommand{\tribManIni}{747,60}
\newcommand{\tribManFim}{1.262,24}
\newcommand{\tribManVar}{68,8}
\newcommand{\transfManIni}{1.855,85}
\newcommand{\transfManFim}{2.739,52}
\newcommand{\transfManVar}{47,6}
\newcommand{\despManIni}{3.382,64}
\newcommand{\despManFim}{4.647,34}
\newcommand{\despManVar}{37,4}
\newcommand{\recManIni}{3.524,47}
\newcommand{\recManFim}{4.757,67}
\newcommand{\recManVar}{35,0}
\newcommand{\fpmManIni}{363,36}
\newcommand{\fpmManFim}{632,22}
\newcommand{\fpmManVar}{74,0}
\newcommand{\pessoalManIni}{1.558,73}
\newcommand{\pessoalManFim}{1.716,33}
\newcommand{\pessoalManVar}{10,1}
\newcommand{\saudeManIni}{611,71}
\newcommand{\saudeManFim}{785,41}
\newcommand{\saudeManVar}{28,4}
\newcommand{\educManIni}{900,07}
\newcommand{\educManFim}{1.248,11}
\newcommand{\educManVar}{38,7}
\newcommand{\investManIni}{388,41}
\newcommand{\investManFim}{492,94}
\newcommand{\investManVar}{26,9}
\newcommand{\fpmShareMan}{19,0}
\newcommand{\icmsShareMan}{46,3}
\newcommand{\fundebShareMan}{27,8}
\newcommand{\susShareMan}{6,9}
\newcommand{\issShareMan}{55,3}
\newcommand{\iptuShareMan}{20,4}
\newcommand{\resPosMan}{6}
\newcommand{\nAnosMan}{10}
\newcommand{\resMinMan}{-219,1}
\newcommand{\resMaxMan}{757,9}
\newcommand{\anoResMinMan}{2022}
\newcommand{\anoResMaxMan}{2021}
\newcommand{\pibpcManIni}{56.787,77}
\newcommand{\pibpcManFim}{64.643,37}
\newcommand{\popOfManVinte}{2.255.903}
\newcommand{\popHaManVinte}{2.038.582}
\newcommand{\popDifMan}{10,7}
\newcommand{\popManFim}{2.303.732}
\newcommand{\anoIniPar}{2017}
\newcommand{\depParMedia}{89,7}
\newcommand{\depParIni}{89,5}
\newcommand{\depParFim}{88,7}
\newcommand{\depParMin}{85,9}
\newcommand{\depParMax}{92,4}
\newcommand{\autParMedia}{5,8}
\newcommand{\autParIni}{6,6}
\newcommand{\autParFim}{6,0}
\newcommand{\tribParIni}{170,05}
\newcommand{\tribParFim}{351,06}
\newcommand{\tribParVar}{106,4}
\newcommand{\transfParIni}{2.324,54}
\newcommand{\transfParFim}{6.210,64}
\newcommand{\transfParVar}{167,2}
\newcommand{\despParIni}{2.745,85}
\newcommand{\despParFim}{6.439,76}
\newcommand{\despParVar}{134,5}
\newcommand{\recParIni}{2.597,80}
\newcommand{\recParFim}{7.005,42}
\newcommand{\recParVar}{169,7}
\newcommand{\fpmParIni}{559,30}
\newcommand{\fpmParFim}{1.250,87}
\newcommand{\fpmParVar}{123,6}
\newcommand{\pessoalParIni}{1.678,42}
\newcommand{\pessoalParFim}{2.764,87}
\newcommand{\pessoalParVar}{64,7}
\newcommand{\saudeParIni}{633,35}
\newcommand{\saudeParFim}{1.649,45}
\newcommand{\saudeParVar}{160,4}
\newcommand{\educParIni}{1.207,60}
\newcommand{\educParFim}{2.242,11}
\newcommand{\educParVar}{85,7}
\newcommand{\investParIni}{44,40}
\newcommand{\investParFim}{1.173,36}
\newcommand{\investParVar}{2.542,7}
\newcommand{\fpmSharePar}{21,1}
\newcommand{\icmsSharePar}{14,3}
\newcommand{\fundebSharePar}{34,2}
\newcommand{\susSharePar}{19,0}
\newcommand{\issSharePar}{36,6}
\newcommand{\iptuSharePar}{6,6}
\newcommand{\resPosPar}{5}
\newcommand{\nAnosPar}{9}
\newcommand{\resMinPar}{-18,9}
\newcommand{\resMaxPar}{57,6}
\newcommand{\anoResMinPar}{2022}
\newcommand{\anoResMaxPar}{2025}
\newcommand{\pibpcParIni}{16.657,63}
\newcommand{\pibpcParFim}{18.343,60}
\newcommand{\popOfParVinte}{116.439}
\newcommand{\popHaParVinte}{96.870}
\newcommand{\popDifPar}{20,2}
\newcommand{\popParFim}{101.855}
\newcommand{\esDepCA}{58,9}
\newcommand{\esDepCAIni}{58,6}
\newcommand{\esDepCAFim}{56,2}
\newcommand{\esAutCA}{35,4}
\newcommand{\esBnoveCAIni}{-0,8}
\newcommand{\esBnoveCAFim}{-0,4}
\newcommand{\esBnoveCAMedia}{-0,5}
\newcommand{\esGDCAIni}{24,7}
\newcommand{\esGDCAFim}{20,2}
\newcommand{\esTECAIni}{15,1}
\newcommand{\esTECAFim}{16,9}
\newcommand{\esDepLoc}{44,3}
\newcommand{\esDepLocIni}{42,2}
\newcommand{\esDepLocFim}{47,9}
\newcommand{\esAutLoc}{47,4}
\newcommand{\esBnoveLocIni}{0,6}
\newcommand{\esBnoveLocFim}{0,3}
\newcommand{\esBnoveLocMedia}{0,3}
\newcommand{\esGDLocIni}{2,9}
\newcommand{\esGDLocFim}{1,2}
\newcommand{\esTELocIni}{5,8}
\newcommand{\esTELocFim}{6,1}
\newcommand{\esDepCen}{11,0}
\newcommand{\esDepCenIni}{8,1}
\newcommand{\esDepCenFim}{12,7}
\newcommand{\esAutCen}{78,5}
\newcommand{\esBnoveCenIni}{-2,5}
\newcommand{\esBnoveCenFim}{-2,0}
\newcommand{\esBnoveCenMedia}{-3,1}
\newcommand{\esGDCenIni}{89,9}
\newcommand{\esGDCenFim}{92,5}
\newcommand{\esTECenIni}{19,6}
\newcommand{\esTECenFim}{22,2}
\newcommand{\esDepGer}{4,9}
\newcommand{\esDepGerIni}{3,9}
\newcommand{\esDepGerFim}{5,3}
\newcommand{\esAutGer}{57,0}
\newcommand{\esBnoveGerIni}{-4,2}
\newcommand{\esBnoveGerFim}{-2,4}
\newcommand{\esBnoveGerMedia}{-4,3}
\newcommand{\esGDGerIni}{102,0}
\newcommand{\esGDGerFim}{100,5}
\newcommand{\esTEGerIni}{42,1}
\newcommand{\esTEGerFim}{45,3}
\newcommand{\esGDvinte}{119,3}
\newcommand{\esBnoveVinte}{-9,9}
\newcommand{\esSupLocAnos}{9}
\newcommand{\bConstUm}{-1,301}
\newcommand{\seConstUm}{2,140}
\newcommand{\pConstUm}{0,555}
\newcommand{\bTrUm}{0,899}
\newcommand{\seTrUm}{0,119}
\newcommand{\pTrUm}{< 0{,}001}
\newcommand{\bPibUm}{0,051}
\newcommand{\sePibUm}{0,039}
\newcommand{\pPibUm}{0,218}
\newcommand{\bParUm}{1,077}
\newcommand{\seParUm}{1,661}
\newcommand{\pParUm}{0,530}
\newcommand{\nUm}{15}
\newcommand{\rdoisUm}{0,870}
\newcommand{\bConstDois}{3,514}
\newcommand{\seConstDois}{7,008}
\newcommand{\pConstDois}{0,627}
\newcommand{\bTrDois}{0,550}
\newcommand{\seTrDois}{0,418}
\newcommand{\pTrDois}{0,218}
\newcommand{\bPibDois}{-0,024}
\newcommand{\sePibDois}{0,113}
\newcommand{\pPibDois}{0,836}
\newcommand{\bParDois}{-1,833}
\newcommand{\seParDois}{4,463}
\newcommand{\pParDois}{0,690}
\newcommand{\bTendDois}{0,145}
\newcommand{\seTendDois}{0,178}
\newcommand{\pTendDois}{0,434}
\newcommand{\nDois}{15}
\newcommand{\rdoisDois}{0,877}
\newcommand{\bConstTres}{-1,346}
\newcommand{\seConstTres}{2,202}
\newcommand{\pConstTres}{0,555}
\newcommand{\bFpmTres}{1,064}
\newcommand{\seFpmTres}{1,040}
\newcommand{\pFpmTres}{0,330}
\newcommand{\bOutTres}{0,847}
\newcommand{\seOutTres}{0,138}
\newcommand{\pOutTres}{< 0{,}001}
\newcommand{\bPibTres}{0,052}
\newcommand{\sePibTres}{0,042}
\newcommand{\pPibTres}{0,240}
\newcommand{\bParTres}{1,125}
\newcommand{\seParTres}{1,846}
\newcommand{\pParTres}{0,556}
\newcommand{\nTres}{15}
\newcommand{\rdoisTres}{0,858}
\newcommand{\bConstQuatro}{1,293}
\newcommand{\seConstQuatro}{7,038}
\newcommand{\pConstQuatro}{0,858}
\newcommand{\bTrQuatro}{0,839}
\newcommand{\seTrQuatro}{0,172}
\newcommand{\pTrQuatro}{< 0{,}001}
\newcommand{\bPibQuatro}{0,042}
\newcommand{\sePibQuatro}{0,745}
\newcommand{\pPibQuatro}{0,956}
\newcommand{\bParQuatro}{-0,300}
\newcommand{\seParQuatro}{0,955}
\newcommand{\pParQuatro}{0,759}
\newcommand{\nQuatro}{15}
\newcommand{\rdoisQuatro}{0,878}
\newcommand{\bConstCinco}{0,143}
\newcommand{\seConstCinco}{0,179}
\newcommand{\pConstCinco}{0,444}
\newcommand{\bTrCinco}{0,515}
\newcommand{\seTrCinco}{0,557}
\newcommand{\pTrCinco}{0,376}
\newcommand{\bPibCinco}{-0,038}
\newcommand{\sePibCinco}{0,129}
\newcommand{\pPibCinco}{0,775}
\newcommand{\nCinco}{13}
\newcommand{\rdoisCinco}{0,011}
\newcommand{\bConstSeis}{-0,612}
\newcommand{\seConstSeis}{0,919}
\newcommand{\pConstSeis}{0,522}
\newcommand{\bTrSeis}{1,040}
\newcommand{\seTrSeis}{0,275}
\newcommand{\pTrSeis}{0,004}
\newcommand{\bPibSeis}{0,034}
\newcommand{\sePibSeis}{0,024}
\newcommand{\pPibSeis}{0,187}
\newcommand{\bParSeis}{0,303}
\newcommand{\seParSeis}{1,052}
\newcommand{\pParSeis}{0,780}
\newcommand{\nSeis}{13}
\newcommand{\rdoisSeis}{0,909}
\newcommand{\bConstSete}{-1,008}
\newcommand{\seConstSete}{0,492}
\newcommand{\pConstSete}{0,065}
\newcommand{\bTrSete}{0,034}
\newcommand{\seTrSete}{0,025}
\newcommand{\pTrSete}{0,204}
\newcommand{\bPibSete}{0,031}
\newcommand{\sePibSete}{0,009}
\newcommand{\pPibSete}{0,005}
\newcommand{\bParSete}{0,556}
\newcommand{\seParSete}{0,386}
\newcommand{\pParSete}{0,178}
\newcommand{\nSete}{15}
\newcommand{\rdoisSete}{0,985}
\newcommand{\bConstOito}{1,440}
\newcommand{\seConstOito}{0,366}
\newcommand{\pConstOito}{0,001}
\newcommand{\bTrOito}{1,076}
\newcommand{\seTrOito}{0,149}
\newcommand{\pTrOito}{< 0{,}001}
\newcommand{\bParOito}{-1,330}
\newcommand{\seParOito}{0,168}
\newcommand{\pParOito}{< 0{,}001}
\newcommand{\nOito}{19}
\newcommand{\rdoisOito}{0,914}
\newcommand{\bConstNove}{-4,063}
\newcommand{\seConstNove}{3,646}
\newcommand{\pConstNove}{0,316}
\newcommand{\bTrNove}{0,051}
\newcommand{\seTrNove}{1,021}
\newcommand{\pTrNove}{0,962}
\newcommand{\bPibNove}{0,129}
\newcommand{\sePibNove}{0,096}
\newcommand{\pPibNove}{0,238}
\newcommand{\nNove}{8}
\newcommand{\rdoisNove}{0,739}
\newcommand{\bConstDez}{2,080}
\newcommand{\seConstDez}{2,387}
\newcommand{\pConstDez}{0,433}
\newcommand{\bTrDez}{1,054}
\newcommand{\seTrDez}{0,315}
\newcommand{\pTrDez}{0,029}
\newcommand{\bPibDez}{-0,108}
\newcommand{\sePibDez}{0,186}
\newcommand{\pPibDez}{0,591}
\newcommand{\nDez}{7}
\newcommand{\rdoisDez}{0,934}
\newcommand{\wtUm}{5,84}
\newcommand{\wpUm}{< 0{,}001}
\newcommand{\wtDois}{1,75}
\newcommand{\wpDois}{0,110}
\newcommand{\wtTres}{0,96}
\newcommand{\wpTres}{0,362}
\newcommand{\wtQuatro}{0,89}
\newcommand{\wpQuatro}{0,395}
\newcommand{\wtCinco}{1,04}
\newcommand{\wpCinco}{0,323}
\newcommand{\wtSeis}{3,41}
\newcommand{\wpSeis}{0,008}
\newcommand{\wtNove}{-0,07}
\newcommand{\wpNove}{0,947}
\newcommand{\wtDez}{2,35}
\newcommand{\wpDez}{0,079}
\newcommand{\razaoFly}{18}
\newcommand{\vifTr}{2,84}
\newcommand{\vifPib}{129,1}
\newcommand{\vifPar}{144,9}
\newcommand{\bpLM}{2,82}
\newcommand{\bpP}{0,420}
\newcommand{\swW}{0,958}
\newcommand{\swP}{0,657}
\newcommand{\resetF}{0,07}
\newcommand{\resetP}{0,801}
\newcommand{\arRho}{-0,33}
\newcommand{\arP}{0,257}
\newcommand{\dwMan}{2,40}
\newcommand{\dwPar}{2,64}
\newcommand{\cookMax}{0,46}
\newcommand{\cookAno}{2021}
\newcommand{\cookMun}{Manaus}
\newcommand{\propTrMan}{1,53}
\newcommand{\propDeMan}{2016}
\newcommand{\dDespMan}{1.110,83}
\newcommand{\dTransfMan}{727,07}
\newcommand{\dPibMan}{7.855,60}
\newcommand{\propTrPar}{0,88}
\newcommand{\propDePar}{2017}
\newcommand{\dDespPar}{1.838,14}
\newcommand{\dTransfPar}{2.091,90}
\newcommand{\dPibPar}{1.685,97}
\newcommand{\esbTr}{1,189}
\newcommand{\esbPib}{0,055}
\newcommand{\esN}{20}
\newcommand{\espTr}{< 0{,}001}
\newcommand{\espPib}{0,024}
\newcommand{\esWt}{6,47}
\newcommand{\bibDocs}{1.992}
\newcommand{\bibFontes}{1109}
\newcommand{\bibAutores}{3.107}
\newcommand{\bibCresc}{7,26}
\newcommand{\bibCit}{18,33}
\newcommand{\bibCoaut}{1,94}
\newcommand{\bibIntl}{8,333}
\newcommand{\bibSingle}{809}
\newcommand{\bibAnoMin}{1952}
\newcommand{\bibAnoMax}{2025}
\newcommand{\bibAteDoisMil}{280}
\newcommand{\bibDesdeVinte}{704}
\newcommand{\bibPctDesdeVinte}{35,3}
\newcommand{\bibFly}{491}
\newcommand{\bibIlu}{309}
\newcommand{\bibDid}{18}
\newcommand{\bibDidDesdeVinte}{15}
\newcommand{\lotkaBeta}{2,83}
\newcommand{\lotkaR}{0,97}
\newcommand{\lotkaP}{0,16}
\newcommand{\fonteA}{Public Choice}
\newcommand{\fonteNA}{54}
\newcommand{\fonteB}{Public Finance Review}
\newcommand{\fonteNB}{28}
\newcommand{\fonteC}{International Tax and Public Finance}
\newcommand{\fonteNC}{26}
\newcommand{\fonteD}{Publius The Journal of Federalism}
\newcommand{\fonteND}{24}
\newcommand{\fonteE}{Public Budgeting \& Finance}
\newcommand{\fonteNE}{23}
\newcommand{\paisA}{Indonesia}
\newcommand{\paisNA}{551}
\newcommand{\paisB}{Usa}
\newcommand{\paisNB}{532}
\newcommand{\paisC}{Brazil}
\newcommand{\paisNC}{461}
\newcommand{\paisD}{Ukraine}
\newcommand{\paisND}{126}
\newcommand{\paisE}{China}
\newcommand{\paisNE}{110}
\newcommand{\paisF}{United Kindom}
\newcommand{\paisNF}{85}
\newcommand{\paisBR}{461}
\newcommand{\paisES}{65}
\newcommand{\tUmTot}{183}
\newcommand{\tUmFly}{74}
\newcommand{\tUmPct}{40,4}
\newcommand{\tUmDezTot}{46}
\newcommand{\tUmDezPct}{62,2}
\newcommand{\tUmFin}{91}
\newcommand{\tUmFinFly}{39}
\newcommand{\tUmFed}{64}
\newcommand{\tUmFedFly}{25}
\newcommand{\tUmBr}{28}
\newcommand{\tUmBrFly}{10}
\newcommand{\tUmBrPct}{35,7}
\newcommand{\corpTot}{1.712}
\newcommand{\corpFly}{640}
\newcommand{\corpPct}{37,4}
\newcommand{\corpDez}{551}
\newcommand{\corpDezPct}{86,1}
\newcommand{\nDid}{18}
\newcommand{\nDidQuinze}{17}
\newcommand{\nBrasil}{244}
\newcommand{\nBrasilFly}{45}
\newcommand{\flEn}{2.466}
\newcommand{\flPt}{406}
\newcommand{\flUnicos}{2.840}
\newcommand{\flSemAut}{15}
\newcommand{\flNaoPer}{549}
\newcommand{\flFinal}{1.992}
\newcommand{\flDoisMilVinteSeis}{284}

\graphicspath{{../07_resultados/figuras/}{../07_resultados/bibliometria/}}

\titulo{Transferências intergovernamentais, ilusão fiscal e efeito \textit{flypaper}: Manaus e Parintins (AM) diante do federalismo fiscal espanhol, 2016--2025}
\autor{João Pedro Furtado Ximenes da Rocha\\ \small Matrícula 231019943}
\data{Outubro de 2026}

\begin{document}
\selectlanguage{brazil}
\frenchspacing

\begin{center}
{\large\bfseries Transferências intergovernamentais, ilusão fiscal e efeito \textit{flypaper}: Manaus e Parintins (AM) diante do federalismo fiscal espanhol, 2016--2025\par}
\vspace{0.4cm}
{\normalsize João Pedro Furtado Ximenes da Rocha --- Matrícula 231019943\par}
{\small Finanças do Setor Público (FSP) --- Prova P2 --- Outubro de 2026\par}
\end{center}
\vspace{0.3cm}

\begin{resumoumacoluna}
O artigo investiga em que medida as transferências intergovernamentais condicionam o financiamento e o gasto de Manaus e Parintins, no Amazonas, entre 2016 e 2025, e compara esse padrão com a organização fiscal dos níveis subnacionais da Espanha. A base do trabalho anterior foi preservada, conferida e reconstruída a partir de fontes oficiais (Siconfi/DCA, IBGE e Eurostat), com harmonização das três classificações da receita pública vigentes no período, exclusão das operações intraorçamentárias, deflação pelo IPCA e pelo IPCH (base 2025) e uma série populacional compatível com os Censos de 2010 e 2022. A literatura foi mapeada com o \textit{bibliometrix} em \bibDocs{} artigos. Em Parintins, as transferências responderam, em média, por \depParMedia\% da receita total, contra \depManMedia\% em Manaus; Em regressões com efeito fixo de município e \nUm{} observações anuais dos dois municípios, cada real adicional de transferência per capita associou-se a R\$ \bTrUm{} de despesa, ao passo que um real adicional de PIB per capita associou-se a R\$ \bPibUm{}, diferença estatisticamente significativa e compatível com o efeito \textit{flypaper}. O padrão vem sobretudo de Parintins, onde a despesa acompanha as transferências quase real a real, e não resiste ao controle de tendência nem às primeiras diferenças, o que recomenda cautela diante do pequeno número de observações. Não se encontrou associação entre transferências e arrecadação tributária própria, que acompanha o PIB local. Na Espanha, as Comunidades Autônomas e as entidades locais têm autonomia tributária muito maior e dependem menos de transferências do que Parintins (\esDepCA\% e \esDepLoc\% da receita, respectivamente) e operam sob regras de estabilidade orçamentária nacionais e europeias. Os resultados indicam associação, e não causalidade, e reforçam a importância de desenhos de transferência que preservem a responsabilização pelo custo do gasto.

\vspace{\onelineskip}
\noindent\textbf{Palavras-chave}: transferências intergovernamentais; efeito \textit{flypaper}; ilusão fiscal; federalismo fiscal; municípios da Amazônia.
\end{resumoumacoluna}

\begin{otherlanguage}{english}
\begin{resumoumacoluna}
This paper examines how intergovernmental transfers shaped the financing and spending of Manaus and Parintins, in the Brazilian state of Amazonas, from 2016 to 2025, and contrasts this pattern with the fiscal organisation of Spain's subnational tiers. The previous assignment's database was preserved, audited and rebuilt from official sources (Siconfi/DCA, IBGE, Eurostat), harmonising three revenue classification schemes, removing intra-budgetary flows, deflating by IPCA and HICP (2025 prices) and using a census-consistent population series. A bibliometric map built with \textit{bibliometrix} covers \bibDocs{} journal articles. Transfers financed on average \depParMedia\% of Parintins' revenue and \depManMedia\% of Manaus'. In regressions with municipality fixed effects on \nUm{} annual observations of the two municipalities, one additional real of transfers per capita is associated with R\$ \bTrUm{} of spending, against R\$ \bPibUm{} for one additional real of local GDP per capita, a statistically significant gap consistent with the flypaper effect. The pattern is driven by Parintins and does not survive a time trend or first differences, which calls for caution given the small sample. Transfers are not associated with own tax revenue. Spanish regional and local governments are less transfer-dependent and operate under domestic and EU fiscal rules. Results are associations, not causal effects.

\vspace{\onelineskip}
\noindent\textbf{Keywords}: intergovernmental transfers; flypaper effect; fiscal illusion; fiscal federalism; Amazon municipalities.
\end{resumoumacoluna}
\end{otherlanguage}

\section{Introdução}\label{sec:intro}

O federalismo fiscal brasileiro atribui aos municípios uma parcela expressiva da provisão de serviços públicos (educação básica, atenção primária à saúde, mobilidade e zeladoria urbana), mas reserva a eles uma base tributária estreita, concentrada no ISS, no IPTU e no ITBI\footnote{BRASIL. \textit{Constituição da República Federativa do Brasil de 1988}. Disponível em: \url{https://www.planalto.gov.br/ccivil_03/constituicao/constituicao.htm}.}. O hiato entre responsabilidades de gasto e capacidade de arrecadação --- o desequilíbrio fiscal vertical --- é coberto por transferências constitucionais e legais da União e dos estados, como o Fundo de Participação dos Municípios (FPM), a cota-parte do ICMS, o FUNDEB e os repasses fundo a fundo do SUS. A teoria do federalismo fiscal reconhece o papel dessas transferências na correção de externalidades e na equalização de capacidades \cite{oates1999,bird2002}, mas a chamada segunda geração da literatura chama atenção para os incentivos que elas criam: governos que gastam recursos arrecadados por outros tendem a perceber o custo do gasto de modo distorcido, a arrecadar menos e a pressionar por novos repasses \cite{weingast2009,rodden2002,oates2008}.

Dois fenômenos concentram esse debate. O primeiro é o efeito \textit{flypaper}: uma transferência incondicional eleva o gasto local muito mais do que um aumento equivalente da renda dos residentes, ao contrário do que prevê o teorema de equivalência derivado do modelo do eleitor mediano \cite{bradford1971,hines1995}. O segundo é a ilusão fiscal, isto é, a percepção sistematicamente distorcida do custo dos bens públicos pelos contribuintes, que pode ser reforçada quando o financiamento vem de outra esfera de governo \cite{wagner1976,logan1986,dollery1996}. No Brasil, a literatura documenta ambos os fenômenos para o conjunto dos municípios \cite{cossio2001,mattos2011,ferreira2019,ferreira2022}, e a tese de Ferreira (2021)\footnote{FERREIRA, L. O. G. \textit{Constitutional tax transfer in Brazilian municipalities: the fiscal illusion and the flypaper effect}. 2021. Tese (Doutorado) -- Universidade de Brasília, Brasília, 2021.}, referência de partida deste trabalho, associa a magnitude do efeito \textit{flypaper} à ilusão fiscal gerada pelas transferências constitucionais.

A Amazônia oferece um contraste particularmente informativo. Manaus, capital do Amazonas e sede da Zona Franca, concentra a atividade industrial do estado e, por isso, tem base própria de ISS e uma cota-parte de ICMS relativamente elevada \cite{serafico2005}. Parintins, município de porte médio situado em uma ilha do baixo Amazonas e conhecido nacionalmente pelo Festival Folclórico, tem PIB per capita cerca de um quarto do de Manaus e depende quase integralmente de recursos de terceiros. As duas cidades compartilham o mesmo estado, as mesmas regras de repartição e o mesmo tribunal de contas, mas diferem radicalmente em escala e em estrutura produtiva, o que permite observar como a dependência de transferências se manifesta em contextos econômicos distintos.

Este trabalho atualiza e amplia a base construída na prova anterior atribuída a esta matrícula\footnote{VICTOR, C. \textit{Prova P1 de Finanças do Setor Público} (matrícula 231020069). 2026. Repositório GitHub: \url{https://github.com/cauaviictor/P1---FSP---231020069---CAUA-VICTOR}, \textit{commit} 5209bcec.}, que analisava a Espanha, Maceió/AL e Duque de Caxias/RJ no período 2013--2022. Com a redistribuição dos conjuntos feita pelo professor, os municípios passaram a ser Manaus e Parintins, mantendo-se a Espanha como país de comparação. A pergunta de pesquisa é: \textbf{em que medida as transferências intergovernamentais condicionam o financiamento e o gasto de Manaus e Parintins entre 2016 e 2025, e esse padrão é compatível com o efeito \textit{flypaper} e com a hipótese de ilusão fiscal, quando confrontado com a organização fiscal subnacional da Espanha?}

O objetivo geral é caracterizar a dependência de transferências dos dois municípios e estimar, com os dados anuais desses dois municípios, a associação entre transferências, renda local e gasto. Os objetivos específicos são: (i) reconstruir, documentar e corrigir a base de dados do trabalho anterior a partir de fontes oficiais, com valores reais e per capita; (ii) mapear a literatura sobre transferências, efeito \textit{flypaper} e ilusão fiscal por meio de um estudo bibliométrico; (iii) descrever a evolução das receitas, das despesas e dos indicadores de autonomia de Manaus e Parintins; (iv) testar se o gasto municipal responde mais às transferências do que à renda local e se as transferências se associam a menor esforço tributário; e (v) contrastar esses padrões com os das Comunidades Autônomas e das entidades locais espanholas.

A justificativa é dupla. Do ponto de vista substantivo, a Reforma Tributária aprovada pela Emenda Constitucional n.º 132/2023 vai substituir o ISS e o ICMS pelo IBS ao longo da próxima década\footnote{BRASIL. Emenda Constitucional n.º 132, de 20 de dezembro de 2023 (Reforma Tributária). Disponível em: \url{https://www.planalto.gov.br/ccivil_03/constituicao/emendas/emc/emc132.htm}.}, alterando justamente as bases que distinguem Manaus de Parintins; uma linha de base bem documentada da dependência atual é condição para avaliar essa transição. Do ponto de vista metodológico, os trabalhos anteriores da disciplina utilizaram planilhas extraídas manualmente, sem registro das classificações contábeis, dos deflatores e das quebras de série; a reconstrução reprodutível aqui proposta permite conferir cada número a partir das fontes. A contribuição do estudo está em analisar em detalhe dois municípios amazônicos raramente estudados, tratar explicitamente as mudanças do plano de contas e da base populacional e situar os resultados em perspectiva internacional com dados do Sistema Europeu de Contas.

A relação entre a análise municipal e a análise do país não pretende comparar Manaus com a Espanha como unidades equivalentes. A Espanha é um Estado unitário fortemente descentralizado, organizado em Comunidades Autônomas e entidades locais\footnote{ESPAÑA. \textit{Constitución Española de 1978}. Disponível em: \url{https://www.boe.es/eli/es/c/1978/12/27/(1)/con}.}, cujo financiamento subnacional combina tributos próprios, impostos cedidos e transferências de nivelamento \cite{zabalza2011,toboso2010}. O país serve, assim, como referência institucional: permite observar como níveis subnacionais com maior autonomia tributária e sujeitos a regras fiscais nacionais e europeias se comportam em termos de dependência, resultado e endividamento. A comparação é feita em proporções da receita e do PIB, que neutralizam as diferenças de moeda e escala, e por subsetor das administrações públicas, aproximando o nível de governo analisado.

O artigo está organizado em quatro seções além desta introdução. A Seção~\ref{sec:referencial} apresenta o estudo bibliométrico e discute a literatura. A Seção~\ref{sec:metodo} descreve o delineamento, as fontes, as variáveis, os tratamentos e as técnicas de análise. A Seção~\ref{sec:resultados} apresenta as mudanças em relação ao trabalho anterior, as estatísticas descritivas, as estimações e as limitações. A Seção~\ref{sec:conclusao} conclui.

\section{Referencial teórico}\label{sec:referencial}

\subsection{Estudo bibliométrico}\label{sec:biblio}

A estimação dos efeitos das transferências intergovernamentais sobre o gasto e a arrecadação dos governos subnacionais está no centro da literatura de federalismo fiscal. As transferências financiam boa parte da provisão local de serviços em países federais e unitários descentralizados, e a forma como os governos receptores as utilizam --- gastando-as, reduzindo tributos ou substituindo esforço próprio --- é decisiva para o desenho dos sistemas de repartição. Duas regularidades organizam esse debate: o efeito \textit{flypaper}, segundo o qual as transferências elevam o gasto muito mais do que um aumento equivalente da renda local, e a ilusão fiscal, segundo a qual o financiamento por terceiros distorce a percepção do custo dos serviços públicos.

A Tabela~\ref{tab:baker}, construída no mesmo formato da Tabela 1 de \citeonline{baker2022}, mostra que, de 2000 a 2025, foram publicados \tUmTot{} artigos sobre transferências intergovernamentais nos quinze periódicos selecionados: \tUmFin{} em periódicos de finanças públicas e escolha pública, \tUmFed{} em periódicos de federalismo, orçamento e governo local e \tUmBr{} em periódicos brasileiros. Entre eles, \tUmFly{} (\tUmPct\%) tratam do efeito \textit{flypaper} ou da ilusão fiscal, e \tUmDezTot{} desses \tUmFly{} (\tUmDezPct\%) foram publicados a partir de 2010. No corpus completo do período, com \corpTot{} artigos em todos os periódicos, a participação é semelhante (\corpFly{} artigos, ou \corpPct\%), e \corpDezPct\% deles são posteriores a 2010. A participação dos dois temas é alta em todos os grupos, inclusive nos periódicos brasileiros (\tUmBrPct\%), o que mostra que o efeito \textit{flypaper} e a ilusão fiscal não são tópicos marginais, mas um eixo central da pesquisa sobre transferências.

\begin{table}[htbp]
\centering
\caption{Artigos sobre transferências intergovernamentais e sobre efeito \textit{flypaper} ou ilusão fiscal, por periódico: 2000--2025}
\label{tab:baker}
\small
\begin{tabular}{lccc}
\toprule
 & (1) & (2) & (3)\\
 & Transferências & \textit{Flypaper}/ilusão fiscal & (2)/(1) (\%)\\
\midrule
\textit{Public Finance Review} & 26 & 13 & 50,0\%\\
\textit{International Tax and Public Finance} & 24 & 5 & 20,8\%\\
\textit{Public Choice} & 22 & 14 & 63,6\%\\
\textit{Journal of Public Economics} & 12 & 3 & 25,0\%\\
\textit{National Tax Journal} & 7 & 4 & 57,1\%\\
\textbf{Finanças públicas e escolha pública} & \textbf{91} & \textbf{39} & \textbf{42,9\%}\\
\midrule
\textit{Public Budgeting \& Finance} & 17 & 10 & 58,8\%\\
\textit{Publius: The Journal of Federalism} & 15 & 3 & 20,0\%\\
\textit{Local Government Studies} & 12 & 7 & 58,3\%\\
\textit{Public Finance and Management} & 13 & 4 & 30,8\%\\
\textit{State and Local Government Review} & 7 & 1 & 14,3\%\\
\textbf{Federalismo, orçamento e governo local} & \textbf{64} & \textbf{25} & \textbf{39,1\%}\\
\midrule
\textit{Economia Aplicada} & 7 & 3 & 42,9\%\\
\textit{Revista de Administração Pública} & 6 & 0 & 0,0\%\\
\textit{Revista Econômica do Nordeste} & 6 & 1 & 16,7\%\\
\textit{Administração Pública e Gestão Social} & 5 & 4 & 80,0\%\\
\textit{Estudos Econômicos} & 4 & 2 & 50,0\%\\
\textbf{Periódicos brasileiros} & \textbf{28} & \textbf{10} & \textbf{35,7\%}\\
\midrule
\textbf{Total dos 15 periódicos} & \textbf{183} & \textbf{74} & \textbf{40,4\%}\\
Corpus completo (todos os periódicos) & 1.712 & 640 & 37,4\%\\
\bottomrule
\end{tabular}
\par\smallskip
\begin{minipage}{0.97\textwidth}\scriptsize
Nota: a Tabela~\ref{tab:baker} segue o formato da Tabela 1 de \citeonline{baker2022}. Resume o número de artigos publicados entre 2000 e 2025 em cinco periódicos de finanças públicas e escolha pública, cinco periódicos de federalismo, orçamento e governo local (os mais produtivos do corpus em cada grupo) e cinco periódicos editados no Brasil. O corpus foi obtido no OpenAlex com busca em título e resumo pelos termos ``intergovernmental transfers'', ``intergovernmental grants'', ``flypaper effect'' e ``fiscal illusion'' e pelos equivalentes em português e espanhol (``transferências intergovernamentais'', ``transferências constitucionais'', ``efeito flypaper'', ``ilusão fiscal'', ``transferencias intergubernamentales'', ``ilusión fiscal''), mantendo apenas artigos com autoria identificada e fonte primária classificada como periódico. A coluna (1) informa o total de artigos do corpus em cada periódico; a coluna (2), quantos deles tratam do efeito \textit{flypaper} ou da ilusão fiscal (termo presente no título, no resumo ou nas palavras-chave); a coluna (3), a participação de (2) em (1). Contagens produzidas com o pacote \textit{bibliometrix} \cite{aria2017} (\texttt{06\_codigos\_e\_planilhas/05\_bibliometria.R}).
\end{minipage}
\end{table}


A difusão desses estudos reflete a percepção de que o efeito \textit{flypaper} é uma das anomalias mais robustas da economia do setor público \cite{hines1995}. Ainda assim, avanços recentes da econometria aplicada mostram que estratégias de identificação muito usadas podem produzir estimativas enganosas. \citeonline{goodmanbacon2021} demonstra que o estimador de efeitos fixos bidirecionais com adoção escalonada é uma média ponderada de comparações entre pares de grupos, algumas delas usando unidades já tratadas como controle; \citeonline{dechaisemartin2020} mostram que os pesos dessas comparações podem ser negativos quando os efeitos são heterogêneos; \citeonline{callaway2021} propõem estimadores por coorte que contornam o problema; e \citeonline{baker2022} mostram que esses vieses podem inverter conclusões publicadas. No corpus analisado, \nDid{} artigos mencionam diferenças-em-diferenças no título, no resumo ou nas palavras-chave, \nDidQuinze{} deles a partir de 2015, o que indica que a literatura sobre transferências começa a adotar esses desenhos. Como as transferências municipais brasileiras não constituem um tratamento discreto adotado em datas distintas, este trabalho não usa diferenças-em-diferenças: estima regressões com efeito fixo de município para Manaus e Parintins e interpreta os coeficientes como associações condicionais (Seção~\ref{sec:metodo}).

O corpus foi construído no OpenAlex\footnote{OPENALEX. API de consulta bibliográfica (\textit{works}). Disponível em: \url{https://api.openalex.org}. Acesso em: 21 out. 2026.} e processado com o \textit{bibliometrix} \cite{aria2017}, por meio do conversor para a API do OpenAlex e das funções \texttt{biblioAnalysis}, \texttt{lotka}, \texttt{biblioNetwork} e \texttt{thematicMap}; o mesmo arquivo pode ser aberto na interface \texttt{biblioshiny()}. As duas consultas (em inglês e em português/espanhol) retornaram \flEn{} e \flPt{} registros, ou \flUnicos{} registros únicos; foram excluídos \flSemAut{} sem autoria, \flNaoPer{} publicados fora de periódicos e \flDoisMilVinteSeis{} de 2026, ano incompleto, restando \flFinal{} artigos (\bibAnoMin--\bibAnoMax). A produção cresce à taxa média de \bibCresc\% ao ano (Figura~\ref{fig:biblio}), \bibPctDesdeVinte\% dos artigos são posteriores a 2019, e \nBrasil{} artigos tratam do Brasil, dos quais \nBrasilFly{} abordam o efeito \textit{flypaper} ou a ilusão fiscal. A distribuição da produtividade dos autores é compatível com a lei de Lotka ($\beta=\lotkaBeta$, $R^2=\lotkaR$, $p=\lotkaP$ no teste de Kolmogorov-Smirnov), e a coautoria internacional é baixa (\bibIntl\% dos artigos). As informações gerais do corpus, a rede de coocorrência de palavras-chave e o mapa temático estão no Apêndice A.

\begin{figure}[htbp]
\centering
\caption{Produção anual de artigos sobre transferências intergovernamentais, efeito \textit{flypaper} e ilusão fiscal, 1970--\bibAnoMax}\label{fig:biblio}
\includegraphics[width=0.88\textwidth]{fig5_bibliometria_producao_anual.png}
\par\smallskip
\begin{minipage}{0.88\textwidth}\scriptsize
Fonte: elaboração própria com \textit{bibliometrix} e OpenAlex. As linhas indicam artigos cujo título, resumo ou palavras-chave mencionam o termo; artigos antigos frequentemente não têm resumo indexado, o que subestima as linhas nos primeiros anos.
\end{minipage}
\end{figure}

\subsection{Federalismo fiscal e o papel das transferências}

A teoria clássica parte da ideia de que a provisão descentralizada permite ajustar a oferta de bens públicos às preferências locais, seja pela mobilidade dos cidadãos entre jurisdições \cite{tiebout1956}, seja pela vantagem informacional dos governos próximos à população \cite{oates1999}. Nesse arcabouço, as transferências cumprem funções bem definidas: corrigir externalidades interjurisdicionais, equalizar capacidades fiscais e cobrir o desequilíbrio vertical decorrente da centralização das bases tributárias mais eficientes \cite{bird2002,bahl2008}. A segunda geração do federalismo fiscal, sintetizada por \citeonline{weingast2009} e \citeonline{oates2008}, incorporou os incentivos políticos e fiscais: governos subnacionais que dependem de transferências têm menos estímulo para promover a base econômica local e mais incentivo para transferir custos ao governo central. A evidência comparada de \citeonline{rodden2002} mostra que déficits persistentes surgem quando os entes subnacionais dependem de transferências gerais e têm liberdade para se endividar, e \citeonline{rodden2003} conclui que a descentralização só contém o crescimento do governo quando é financiada por tributos próprios, e não por repasses. Na América Latina, \citeonline{stein1999} associa o desequilíbrio vertical ao problema do fundo comum, que amplia o tamanho do governo, e os modelos de economia política do orçamento explicam por que a dispersão do custo e a concentração do benefício geram viés deficitário \cite{alesina1995,vonhagen1995}.

A literatura também mostra que a alocação das transferências não é neutra. Governos centrais usam repasses discricionários de forma tática para favorecer eleitores decisivos \cite{johansson2003} e prefeitos alinhados politicamente, como documentado para a Espanha \cite{soleolle2008,simon2013}, para o Brasil \cite{brollo2012} e em modelos de sinalização com evidência italiana \cite{bracco2015}. Mesmo as transferências constitucionais, distribuídas por fórmula, produzem efeitos que vão além do volume de gasto: no Brasil, os saltos do FPM em faixas populacionais elevaram a corrupção observada e reduziram a escolaridade média dos candidatos a prefeito \cite{brollo2013}, ao mesmo tempo em que aumentaram o gasto, a escolarização e a alfabetização e reduziram a pobreza, sem deslocar a receita própria \cite{litschig2013}, com ganhos que persistem no longo prazo \cite{carvalhofilho2022}. Auditorias federais indicam ainda que a responsabilização eleitoral reduz o desvio de recursos transferidos \cite{ferraz2011}. Esses resultados sugerem que o efeito das transferências depende do ambiente institucional em que são aplicadas, ponto retomado por \citeonline{lago2024} em ampla revisão recente.

\subsection{O efeito \textit{flypaper}: teoria, evidência e explicações}

O ponto de partida é o teorema de equivalência de \citeonline{bradford1971}: se o governo local apenas agrega as preferências do eleitor mediano, uma transferência incondicional equivale a um aumento da renda dos residentes e deveria ser gasta na mesma proporção que a renda privada, isto é, alguns centavos por real. A evidência empírica contrariou sistematicamente essa previsão. Na síntese de \citeonline{hines1995}, enquanto um real adicional de renda eleva o gasto estadual americano em 5 a 10 centavos, um real de transferência eleva-o em valores que chegam a um real --- daí a imagem de que o dinheiro ``gruda onde cai''. Estudos para municípios italianos \cite{gennari2014}, mexicanos \cite{cardenas2011}, quebequenses \cite{mehiriz2014}, alemães \cite{baskaran2015} e do estado de Wisconsin \cite{deller2005} reproduzem o fenômeno, e ele aparece também no gasto regional com saúde na Itália, acompanhado de ajuste lento \cite{levaggi2003}; no caso das transferências educacionais, a resposta depende do tipo de repasse e das regras de contrapartida \cite{fisher2000}. Mesmo desenhos quase-experimentais confirmam o efeito: a regra descontínua de transferências suecas elevou o gasto sem reduzir impostos \cite{dahlberg2008}, e os estados americanos gastaram mais de um dólar por dólar recebido do programa rodoviário de estímulo de 2009 \cite{leduc2017}. A evidência não é unânime: \citeonline{knight2002} mostra que, ao tratar a endogeneidade das transferências federais com instrumentos políticos, o efeito desaparece e há deslocamento significativo do gasto estadual; \citeonline{zampelli1986} encontra pouca evidência do efeito quando se considera a fungibilidade de transferências condicionais; e \citeonline{gamkhar1996} não encontram assimetria entre aumentos e cortes, ao passo que \citeonline{heyndels2001} identifica assimetria em municípios flamengos, com gasto que sobe quando as transferências crescem e impostos que sobem quando elas caem.

As explicações se dividem em quatro famílias. A primeira atribui o efeito a problemas de especificação e de estimação: \citeonline{fisher1982} e \citeonline{hamilton1983} questionam a comparabilidade entre renda e transferência nos modelos usuais, \citeonline{moffitt1984} mostra que restrições orçamentárias não lineares de programas com contrapartida distorcem as estimativas, \citeonline{megdal1987} demonstra, por simulação, que o MQO superestima a propensão a gastar quando a transferência é construída endogenamente, e \citeonline{becker1996} argumenta que parte do que se chama ilusão fiscal é artefato de especificação. A segunda enfatiza a assimetria de informação e o poder de agenda dos burocratas, que conhecem o tamanho da transferência e a utilizam para expandir o orçamento \cite{romer1980,filimon1982}; \citeonline{strumpf1998} deriva desse arcabouço um índice preditivo da magnitude do efeito, e \citeonline{wyckoff1991} confronta empiricamente as explicações concorrentes. A terceira explica o efeito pelo custo marginal dos recursos públicos: como os tributos locais geram peso morto, um real transferido custa menos ao eleitor do que um real arrecadado, o que torna racional gastar mais da transferência \cite{hamilton1986,dahlby2011}; o custo elevado da arrecadação local amplia a resposta às transferências em municípios peruanos \cite{aragon2013}. A quarta família, comportamental, associa o efeito à contabilidade mental --- recursos ``do grupo'' são tratados como distintos dos recursos individuais, como mostram experimentos de laboratório \cite{becker2020} --- e à incerteza: quando as transferências são menos voláteis ou menos correlacionadas com a renda privada, a poupança precaucional responde de modo diferente aos dois choques \cite{vegh2015}. Há ainda interpretações institucionais, segundo as quais as instituições políticas moderam a resposta às transferências e ajudam a explicar a evidência mista \cite{kjaergaard2015}, sendo o efeito menor onde a competição política e a eficiência arrecadatória são maiores \cite{rios2021}, e a constatação de que as transferências de equalização se capitalizam no preço dos imóveis, o que altera a leitura do efeito \cite{allers2016}. Revisões da literatura mantêm o diagnóstico de que o efeito existe, mas sua magnitude depende do desenho das transferências e da estratégia de identificação \cite{bailey1998,lago2024}.

\subsection{Ilusão fiscal}

A hipótese de ilusão fiscal sustenta que os contribuintes subestimam o custo dos serviços públicos quando a estrutura de receitas é complexa, pouco visível ou financiada por terceiros \cite{wagner1976}. O estudo pioneiro de \citeonline{pommerehne1978} para cidades suíças encontrou gastos maiores onde a carga tributária era menos perceptível, resultado que estimulou testes para a complexidade tributária \cite{misiolek1988,heyndels1995}, para municípios flamengos \cite{heyndels1994}, para o Reino Unido \cite{gemmell1999} e para a dívida pública local \cite{banzhaf2013}. \citeonline{dollery1996} organizam a literatura em cinco hipóteses --- complexidade da receita, elasticidade da receita, efeito \textit{flypaper}, ilusão do inquilino e ilusão da dívida ---, das quais a terceira é a que interessa diretamente a este trabalho. Nela, as transferências reduzem o preço percebido do bem público local \cite{logan1986,turnbull1992,turnbull1998}, e o contribuinte que não distingue os recursos recebidos de outras esferas tende a apoiar orçamentos maiores; \citeonline{logan1986} mostra que, simetricamente, o governo que concede a transferência parece mais caro, efeito confirmado com dados de pesquisa por \citeonline{hewitt1986}. A ilusão é maior quando a transferência vem de uma esfera mais distante \cite{grossman1990}, e \citeonline{mitias2001} identificam que o gasto dos condados americanos responde à ilusão de transferência, e não à ilusão tributária.

A evidência é heterogênea e alimenta um debate sobre a própria existência da ilusão. Experimentos de laboratório indicam que tributos menos visíveis são subestimados pelos participantes e distorcem suas escolhas de gasto \cite{sausgruber2005}, e \citeonline{baekgaard2016} mostram que ela decorre menos da falta de informação do que da falta de atenção. Índices construídos por modelos de múltiplos indicadores mostram grande variação entre países \cite{dellanno2012} e na União Europeia \cite{dellanno2014}. Para governos locais, o fenômeno aparece em municípios australianos \cite{worthington1999} e sul-africanos, em que a forte dependência de transferências é vista como fonte de ilusão \cite{amusa2008}. No Brasil, há evidência de ilusão fiscal associada a transferências e déficits nos estados \cite{prado2020} e na demanda por gasto federal \cite{silva2015}; nos municípios, a captura de transferências por grupos de interesse cresce com a ilusão fiscal e com o sobrefinanciamento de alguns entes \cite{mendes2005}, e a natureza das transferências relaciona-se com o cumprimento das metas fiscais \cite{araujo2020}.

\subsection{Evidências para os municípios brasileiros}

O estudo de \citeonline{cossio2001} é uma das primeiras evidências do efeito \textit{flypaper} no Brasil: as transferências constitucionais de natureza incondicional têm efeito mais expansivo sobre as despesas municipais do que aumentos equivalentes na renda, com transbordamentos espaciais entre municípios vizinhos. \citeonline{mattos2011} reinterpretam o fenômeno como perda de eficiência arrecadatória, e não apenas como expansão do gasto, ao mostrar que municípios que recebem mais transferências exploram menos sua base tributária; \citeonline{sakurai2013} encontra respostas assimétricas a aumentos e reduções de transferências, com recomposição entre gasto corrente e investimento; e \citeonline{nojosa2018} mostram que o efeito é mais intenso onde o prefeito tem maior força política na câmara. \citeonline{ferreira2019} confirmam o efeito em municípios brasileiros com instrumentos de alinhamento político, intensificado quando prefeito, governador e presidente são do mesmo partido, e \citeonline{ferreira2022} constroem um índice do efeito para estados e municípios entre 2000 e 2018, documentando sua relevância generalizada. A tese de Ferreira (2021)\footnote{FERREIRA, L. O. G. \textit{Constitutional tax transfer in Brazilian municipalities: the fiscal illusion and the flypaper effect}. 2021. Tese (Doutorado) -- Universidade de Brasília, Brasília, 2021.} articula esses resultados à ilusão fiscal produzida pelas transferências constitucionais.

Outros estudos qualificam o mecanismo. As transferências em bloco ampliam a dispersão do gasto educacional entre municípios, ao contrário das condicionais \cite{mattos2018}; a regra de gasto mínimo em educação elevou o gasto apenas nos municípios que estavam abaixo do piso \cite{cruz2020}; parte do efeito \textit{flypaper} decorre de transbordamentos das transferências recebidas por municípios vizinhos \cite{castro2021}; as faixas do FPM afetam a composição do gasto e os incentivos fiscais locais \cite{alves2024}; e a descentralização altera a composição da despesa de forma heterogênea ao longo da distribuição \cite{sousa2021}. Receitas de royalties, que funcionam como transferências não vinculadas, reduziram o esforço tributário de municípios produtores de petróleo \cite{postali2015}, o que reforça a hipótese de substituição. Na saúde, o financiamento federal cresceu sobretudo por transferências, que superaram 70\% do gasto federal no setor na década de 2000 \cite{machado2014}. Do ponto de vista do gasto agregado, \citeonline{mello2002} relaciona as fontes de financiamento à composição do gasto local e ao crescimento.

A disciplina fiscal municipal brasileira se apoia na Lei de Responsabilidade Fiscal\footnote{BRASIL. Lei Complementar n.º 101, de 4 de maio de 2000 (Lei de Responsabilidade Fiscal). Disponível em: \url{https://www.planalto.gov.br/ccivil_03/leis/lcp/lcp101.htm}.}, cujos efeitos sobre o endividamento municipal foram documentados para os municípios com mais de 100 mil habitantes \cite{macedo2009} e para o Piauí \cite{linhares2013}. A lei se insere na agenda de administração gerencial que buscou deslocar o foco do controle de processos para o de resultados \cite{bresserpereira1996}\footnote{BRASIL. Presidência da República. Câmara da Reforma do Estado. \textit{Plano Diretor da Reforma do Aparelho do Estado}. Brasília, 1995.}. A literatura bibliométrica nacional sobre transferências, por sua vez, concentra-se nas transferências da União para os municípios e na aplicação dos recursos em saúde e educação \cite{gomes2012}.

\subsection{A experiência espanhola}

A Espanha construiu, a partir de 1978, um Estado autonômico em que as Comunidades Autônomas assumiram saúde, educação e serviços sociais\footnote{ESPAÑA. \textit{Constitución Española de 1978}. Disponível em: \url{https://www.boe.es/eli/es/c/1978/12/27/(1)/con}.}. Até as reformas de 1996 e 2001, as regiões de regime comum eram financiadas majoritariamente por transferências; desde então, passaram a receber parcelas cedidas do IRPF, do IVA e dos impostos especiais, com capacidade normativa limitada \cite{toboso2010}. A reforma de 2009 reorganizou os fundos de garantia e de suficiência \cite{zabalza2011}\footnote{ESPAÑA. Ley 22/2009, de 18 de diciembre, que regula o sistema de financiação das Comunidades Autónomas de regime comum. Disponível em: \url{https://www.boe.es/eli/es/l/2009/12/18/22/con}.}, e a maior autonomia tributária endureceu a restrição orçamentária regional, reduzindo o incentivo a acumular déficits \cite{arespa2022}. Esse resultado converge com a evidência para a OCDE de que a descentralização ``parcial'', financiada por transferências, compromete a disciplina fiscal subnacional, ao passo que a autonomia tributária a fortalece \cite{asatryan2015}. A descentralização também se associa ao tamanho dos governos regionais espanhóis \cite{cantarero2012}, e a distribuição das transferências discricionárias responde a estratégias eleitorais \cite{simon2013}.

No nível local, os municípios espanhóis combinam tributos próprios (IBI, IAE, imposto sobre veículos e sobre construções) com a participação nos tributos do Estado\footnote{ESPAÑA. Real Decreto Legislativo 2/2004, de 5 de marzo, Texto Refundido de la Ley Reguladora de las Haciendas Locales. Disponível em: \url{https://www.boe.es/eli/es/rdlg/2004/03/05/2/con}.}. A fórmula de repartição pondera a população com pesos crescentes, refletindo custos maiores dos municípios grandes \cite{soleolle2005}; a expansão urbana dispersa elevou receitas vinculadas à construção e despesas de infraestrutura \cite{hortasrico2014}; e a situação financeira municipal se explica tanto por fatores políticos quanto socioeconômicos \cite{bastida2009}. Para 2.451 municípios, \citeonline{rios2021} mostram que o efeito \textit{flypaper} é menor onde a eficiência arrecadatória e a força política do governo são maiores, enquanto a fragmentação partidária reduz a estabilidade dos governos locais \cite{carozzi2022}. Desde 2012, a Lei Orgânica de Estabilidade Orçamentária impõe a todas as administrações limites de déficit, dívida e uma regra de gasto\footnote{ESPAÑA. Ley Orgánica 2/2012, de 27 de abril, de Estabilidad Presupuestaria y Sostenibilidad Financiera. Disponível em: \url{https://www.boe.es/eli/es/lo/2012/04/27/2/con}.}, articulada ao Pacto de Estabilidade e Crescimento e ao novo arcabouço de governança econômica europeu\footnote{UNIÃO EUROPEIA. Regulamento (UE) 2024/1263 relativo à coordenação eficaz das políticas económicas e à supervisão orçamental multilateral. \textit{Jornal Oficial da União Europeia}, L 2024/1263.}. O ajuste fiscal das administrações locais depende, em boa medida, de como as transferências reagem a choques de receita e despesa, como mostra \citeonline{buettner2006} para municípios americanos, e a experiência da África Subsaariana sugere que, onde a capacidade fiscal local é limitada, a relação entre transferências e arrecadação própria não segue necessariamente o padrão de deslocamento observado em países de renda alta \cite{masaki2018}.

\subsection{Síntese, lacunas e hipóteses}

A literatura converge em três pontos. O gasto subnacional tende a responder mais às transferências do que à renda local; essa resposta é mais forte quando a base própria é estreita e cara de explorar; e a dependência de transferências enfraquece a disciplina fiscal na ausência de regras ou de autonomia tributária. Divergem, porém, as interpretações sobre se o efeito reflete ilusão dos eleitores, poder de agenda dos burocratas, custo dos recursos ou problemas de identificação. Há duas lacunas a que este trabalho procura responder: os estudos brasileiros raramente analisam municípios da Região Norte, cujas condições de arrecadação e de custo de provisão são muito particulares, e as comparações internacionais pouco exploram dados por subsetor das contas nacionais. Com base nessa discussão, são formuladas as seguintes hipóteses verificáveis:

\begin{description}
\item[H1 (dependência)] Parintins é mais dependente de transferências do que Manaus, e ambos são mais dependentes do que as Comunidades Autônomas e as entidades locais espanholas.
\item[H2 (\textit{flypaper})] Em Manaus e Parintins, a propensão marginal a gastar as transferências supera a propensão a gastar a renda local ($\beta_1>\beta_2$).
\item[H3 (ilusão fiscal e esforço tributário)] Transferências maiores associam-se a menor arrecadação tributária própria per capita ($\gamma_1<0$).
\item[H4 (disciplina fiscal)] Os níveis subnacionais espanhóis, com maior autonomia tributária e regras de estabilidade, apresentam resultados mais próximos do equilíbrio do que a administração central.
\end{description}

\section{Metodologia}\label{sec:metodo}

\subsection{Delineamento, unidades de análise e período}

Trata-se de pesquisa quantitativa, de natureza descritiva e explicativa, com delineamento longitudinal de estudo de casos múltiplos: dois municípios brasileiros (Manaus e Parintins) acompanhados ano a ano e um país de referência (Espanha), analisado por subsetor das administrações públicas. As unidades de análise são, portanto, apenas as três atribuídas no trabalho. Como cada município contribui com no máximo dez observações anuais, a análise privilegia a descrição detalhada das séries e usa as regressões como evidência complementar, com especificações simples, erros-padrão adequados a amostras pequenas e testes de sensibilidade explícitos.

O período é 2016--2025, os dez exercícios mais recentes com Declaração de Contas Anuais (DCA) completa no Siconfi na data da coleta (21 de outubro de 2026). As coberturas diferem entre variáveis e unidades, e essas diferências não foram preenchidas artificialmente: (i) Parintins não tem DCA de 2016 homologada no Siconfi --- a API do Tesouro não retorna registros nem para a DCA nem para o RREO do 6.º bimestre, e o extrato de entregas do ente registra para aquele ano apenas os relatórios de gestão fiscal ---, de modo que sua série começa em 2017; (ii) o PIB municipal do IBGE está disponível até 2023, o que limita as regressões que usam a renda local a 2016--2023 (15 observações: 8 de Manaus e 7 de Parintins); e (iii) para a Espanha, todas as séries do Eurostat estão disponíveis até 2025.

\subsection{Fontes e critérios de seleção dos dados}

Os dados municipais foram obtidos diretamente na API de dados abertos do Siconfi\footnote{SECRETARIA DO TESOURO NACIONAL. API de Dados Abertos do Siconfi -- Declaração de Contas Anuais (DCA), Anexos I-C, I-D e I-E. Disponível em: \url{https://apidatalake.tesouro.gov.br/ords/siconfi/tt/dca}. Acesso em: 21 out. 2026.}, que replica as Declarações de Contas Anuais enviadas pelos municípios. Foram usados os Anexos I-C (receitas orçamentárias por natureza), I-D (despesas por natureza) e I-E (despesas por função) da DCA, que consolidam o exercício completo, em vez do RREO, que pode ser retificado após o fechamento. A população, o PIB municipal e o IPCA vêm das tabelas 6579, 4714, 1378, 5938 e 1737 do Sidra/IBGE\footnote{INSTITUTO BRASILEIRO DE GEOGRAFIA E ESTATÍSTICA. Sistema IBGE de Recuperação Automática (Sidra): tabelas 1737, 6579, 4714, 1378 e 5938. Disponível em: \url{https://servicodados.ibge.gov.br/api/v3/agregados}. Acesso em: 21 out. 2026.}; os dados da Espanha vêm das tabelas \texttt{gov\_10a\_main}, \texttt{gov\_10dd\_edpt1}, \texttt{prc\_hicp\_aind}, \texttt{nama\_10\_gdp} e \texttt{demo\_pjan} do Eurostat\footnote{EUROSTAT. Tabelas gov\_10a\_main, gov\_10dd\_edpt1, prc\_hicp\_aind, nama\_10\_gdp e demo\_pjan. Disponível em: \url{https://ec.europa.eu/eurostat/api/dissemination/statistics/1.0/data/}. Acesso em: 21 out. 2026.}, compiladas segundo o Sistema Europeu de Contas (SEC 2010). Optou-se pelo Eurostat, e não pelo \textit{World Economic Outlook} do FMI utilizado no trabalho anterior, porque o Eurostat fornece os dados por subsetor (administração central, Comunidades Autônomas, entidades locais e seguridade social), indispensáveis para comparar níveis de governo, e porque suas séries são as notificadas no âmbito do procedimento de déficit excessivo. Todas as consultas, com URL, data e parâmetros, estão registradas nos arquivos brutos e no dicionário de variáveis do repositório.

\subsection{Variáveis e formas de mensuração}

O Quadro~\ref{tab:variaveis} apresenta as principais variáveis. O dicionário completo, com conceito, fonte, endereço eletrônico, período, cobertura, unidade, tratamento de ausentes e fórmula de cada variável, está no arquivo \texttt{04\_dicionario\_variaveis/dicionario\_variaveis.xlsx}.

\begin{table}[htbp]
\centering
\caption{Principais variáveis, definições e fontes}\label{tab:variaveis}
\scriptsize
\begin{tabular}{p{3.1cm}p{8.2cm}p{3.2cm}}
\toprule
Variável & Definição e mensuração & Fonte\\
\midrule
Receita total líquida & Receitas correntes + de capital, exceto intraorçamentárias, deduzidas as transferências ao FUNDEB e outras deduções (R\$) & Siconfi, DCA I-C\\
Receita tributária & Impostos, taxas e contribuições de melhoria (código 1.1), inclusive IRRF dos servidores & Siconfi, DCA I-C\\
Arrecadação própria corrente & Receita corrente bruta menos transferências correntes brutas & Siconfi, DCA I-C\\
Transferências líquidas & Transferências correntes (1.7) líquidas das deduções para o FUNDEB + transferências de capital (2.4) & Siconfi, DCA I-C\\
FPM, ICMS, FUNDEB, SUS & Rubricas identificadas pelo nome da conta nas três classificações vigentes (2016--2017, 2018--2021 e 2022--2025) & Siconfi, DCA I-C\\
Despesa empenhada & Despesas exceto intraorçamentárias, por natureza (pessoal, investimentos) e por função (saúde, educação) & Siconfi, DCA I-D e I-E\\
Dependência & Transferências líquidas / receita total líquida & Calculada\\
Autonomia tributária & Receita tributária / receita corrente líquida & Calculada\\
Resultado orçamentário & Receita total líquida menos despesa empenhada & Calculada\\
PIB per capita & PIB a preços correntes / população harmonizada & IBGE, tab. 5938\\
População harmonizada & Interpolação geométrica entre os Censos de 2010 e 2022 (2016--2021); Censo 2022; interpolação 2022--2024 (2023); estimativas pós-censo (2024--2025) & IBGE, tab. 1378, 4714 e 6579\\
Deflatores & IPCA, média anual do número-índice, base 2025 = 1; IPCH (HICP) da Espanha, média anual, base 2025 = 1 & IBGE tab. 1737; Eurostat\\
Espanha (por subsetor) & Receita total (TR), despesa total (TE), saldo (B.9), dívida (GD), impostos (D.2 + D.5 + D.91), transferências recebidas (D.7 + D.9), em \% do PIB e em euros de 2025 per capita & Eurostat\\
\bottomrule
\end{tabular}
\par\smallskip
\begin{minipage}{0.97\textwidth}\scriptsize Fonte: elaboração própria.\end{minipage}
\end{table}

\subsection{Tratamentos e transformações}\label{sec:tratamentos}

Cinco tratamentos merecem destaque, todos implementados em código e documentados no registro de alterações.

\textbf{Harmonização do plano de contas.} A classificação da receita mudou duas vezes no período: até 2017 vigorava a codificação herdada da Portaria Interministerial n.º 163/2001; entre 2018 e 2021, códigos de transição (por exemplo, 1.7.1.8.01.2.0 para o FPM); e, a partir de 2022, a estrutura definitiva (1.7.1.1.51.0.0 para o FPM). Como o mesmo conceito tem códigos diferentes em cada fase, as rubricas detalhadas foram identificadas pelo nome da conta, e não pelo código, com uma regra que soma apenas a conta de nível mais alto quando uma conta e suas subcontas atendem ao mesmo critério, evitando dupla contagem. As contas efetivamente somadas em cada ano estão listadas no arquivo de auditoria.

\textbf{Exclusão das operações intraorçamentárias.} Os agregados de despesa do Anexo I-D incluem as despesas intraorçamentárias (modalidade 91), sobretudo contribuições patronais ao regime próprio de previdência, que também aparecem como receita intraorçamentária. Para que receita e despesa sejam comparáveis, ambas foram tomadas exceto intraorçamentárias. Em Manaus, essa correção retira da despesa total entre R\$ 0,16 bilhão (2017) e R\$ 0,90 bilhão (2020) por ano, em valores nominais; Parintins não registra operações intraorçamentárias no período.

\textbf{Deduções do FUNDEB.} Os municípios contribuem compulsoriamente com 20\% de parte de suas transferências para o FUNDEB e recebem de volta valores redistribuídos dentro do estado\footnote{BRASIL. Lei n.º 14.113, de 25 de dezembro de 2020 (regulamenta o FUNDEB permanente). Disponível em: \url{https://www.planalto.gov.br/ccivil_03/_ato2019-2022/2020/lei/l14113.htm}.}. As receitas foram consideradas líquidas dessas deduções, e o FUNDEB recebido foi mantido entre as transferências, de modo que a dependência mede o que efetivamente financia o município.

\textbf{Deflação.} Os valores monetários brasileiros foram convertidos para reais de 2025 pelo fator $F_t = \bar I_{2025}/\bar I_t$, em que $\bar I_t$ é a média dos doze números-índices mensais do IPCA no ano $t$. Usa-se a média anual porque receitas e despesas são fluxos distribuídos ao longo do ano. Os valores espanhóis foram deflacionados da mesma forma pelo IPCH. O trabalho anterior usava fatores pontuais da Calculadora do Cidadão com base 2022 e atribuía fator diferente de 1 ao próprio ano-base, o que foi corrigido.

\textbf{População e valores per capita.} As estimativas anuais do IBGE para 2016--2021 foram projetadas a partir do Censo de 2010, e o Censo de 2022 revelou que elas superestimavam a população: em 2021, a estimativa oficial de Manaus excedia em \popDifMan\% o valor compatível com os dois censos, e a de Parintins, em \popDifPar\%. Usar as estimativas antigas e o Censo de 2022 na mesma série criaria uma quebra artificial nos valores per capita. Por isso, a série principal é uma população harmonizada, obtida por interpolação geométrica entre os censos (com datas de referência de 1.º de agosto para os censos e de 1.º de julho para as estimativas); para 2023, ano sem publicação oficial, usa-se a interpolação entre o Censo de 2022 e a estimativa de 2024. A série oficial é mantida na base e usada como teste de sensibilidade. Dados ausentes nunca foram substituídos por zero: permanecem como ausentes e as estatísticas usam apenas as observações disponíveis.

\subsection{Técnicas de análise}

A análise descritiva compara a evolução das receitas, das transferências e das despesas reais per capita de Manaus e Parintins entre si e com os subsetores espanhóis, e calcula, para cada município, a variação acumulada da despesa por real de variação das transferências entre o primeiro e o último ano com PIB disponível. Para testar as hipóteses H2 e H3, estima-se por mínimos quadrados ordinários (MQO), com os dados dos dois municípios empilhados (município-ano), o modelo
\begin{equation}\label{eq:fly}
G_{it} = \alpha + \delta P_i + \beta_1 T_{it} + \beta_2 Y_{it} + \varepsilon_{it},
\end{equation}
em que $G_{it}$ é a despesa empenhada real per capita do município $i$ no ano $t$, $T_{it}$ são as transferências líquidas reais per capita, $Y_{it}$ é o PIB real per capita (\textit{proxy} da renda local) e $P_i$ é uma variável indicadora de Parintins, que funciona como efeito fixo de município e absorve as diferenças permanentes de escala, de custo e de estrutura administrativa entre as duas cidades. Os coeficientes são propensões marginais a gastar: o teorema de equivalência prevê $\beta_1=\beta_2$, e o efeito \textit{flypaper} corresponde a $\beta_1>\beta_2$, com sinais esperados $\beta_1>0$ e $\beta_2>0$. A hipótese é testada pela estatística $t$ de $H_0\!:\beta_1-\beta_2=0$. A equação de esforço tributário é
\begin{equation}\label{eq:trib}
R_{it} = \alpha + \delta P_i + \gamma_1 T_{it} + \gamma_2 Y_{it} + u_{it},
\end{equation}
em que $R_{it}$ é a receita tributária real per capita. A hipótese de ilusão fiscal com substituição de esforço prevê $\gamma_1<0$; a ausência de efeito ($\gamma_1=0$) é compatível com transferências que não deslocam a arrecadação, como encontrado por \citeonline{litschig2013}.

Com 15 observações, os erros-padrão usuais e os assintóticos são pouco confiáveis. Por isso, usam-se erros robustos à heterocedasticidade do tipo HC3, os mais conservadores da família, e valores críticos da distribuição $t$ de Student com os graus de liberdade efetivos. Não se usam erros agrupados por município nem os de \citeonline{driscoll1998}, porque ambos exigem muitas unidades (ou muitos períodos) para serem consistentes, e aqui há apenas duas. A sensibilidade é avaliada por oito variações: (i) inclusão de tendência linear comum, que controla o crescimento simultâneo de receitas e despesas; (ii) separação do FPM, transferência incondicional típica, das demais transferências; (iii) especificação em logaritmos, cujos coeficientes são elasticidades; (iv) primeiras diferenças, que eliminam tendências e efeitos fixos; (v) valores per capita calculados com a população oficial do IBGE; (vi) e (vii) estimação separada para cada município; e (viii) exclusão do PIB, o que permite usar todos os anos até 2025 (19 observações). Os diagnósticos incluem fatores de inflação da variância (VIF), o teste de Breusch-Pagan para heterocedasticidade, o de Shapiro-Wilk para normalidade dos resíduos, o RESET de Ramsey para forma funcional, a autocorrelação de primeira ordem dos resíduos dentro de cada município (coeficiente AR(1) e estatística de Durbin-Watson) e a distância de Cook para observações influentes.

Para a Espanha, com dez anos por subsetor, estima-se uma regressão análoga, por MQO com erros robustos e efeito fixo de subsetor, da despesa real per capita das Comunidades Autônomas e das entidades locais sobre as transferências recebidas e o PIB per capita ($n=\esN$), lida apenas como indício descritivo da mesma regularidade.

\subsection{Associação, causalidade e comparabilidade}

Os coeficientes da Equação~\eqref{eq:fly} medem associações condicionais, e não efeitos causais. As transferências não são atribuídas aleatoriamente: o FPM depende da faixa populacional, a cota-parte do ICMS depende do valor adicionado e de critérios estaduais, e os repasses do SUS e do FUNDEB dependem de matrículas e de produção de serviços --- variáveis que também afetam o gasto. A variável indicadora de município remove as diferenças permanentes entre Manaus e Parintins, mas não as variações ao longo do tempo que afetam simultaneamente transferências e despesas, como o crescimento da população atendida ou mudanças nas regras de financiamento da saúde e da educação. Com apenas dois municípios, também não é possível controlar choques comuns por efeitos de ano, o que torna a tendência linear e as primeiras diferenças testes importantes. Além disso, parte das transferências é vinculada a gastos específicos, o que mecanicamente aproxima $\beta_1$ de 1 e torna o teste mais exigente de interpretar \cite{knight2002,megdal1987}. Desenhos quase-experimentais, como as descontinuidades do FPM \cite{brollo2013,litschig2013}, seriam necessários para uma interpretação causal; como discutido na Seção~\ref{sec:biblio}, esses desenhos têm suas próprias exigências \cite{baker2022,goodmanbacon2021}.

A comparação entre municípios e país também tem limites. As moedas, as escalas e os níveis de preços são diferentes, o que se contorna com proporções (da receita e do PIB) e com valores deflacionados por índices nacionais, sem conversão cambial. As contas brasileiras seguem a contabilidade orçamentária e patrimonial do MCASP\footnote{SECRETARIA DO TESOURO NACIONAL. \textit{Manual de Contabilidade Aplicada ao Setor Público (MCASP)}. 10.ª ed. Brasília: STN, 2023.}, de base mista, ao passo que as espanholas seguem o SEC 2010, de competência econômica. A organização fiscal também difere: as entidades locais espanholas têm competências mais restritas que os municípios brasileiros (saúde e educação são regionais), e o SEC registra como imposto regional a parcela cedida do IRPF e do IVA, que no Brasil seria classificada como transferência. Por isso, os municípios brasileiros são comparados com os dois subsetores espanhóis --- as Comunidades Autônomas, pela natureza das competências, e as entidades locais, pelo nível de governo ---, e as conclusões comparativas são qualitativas.

\subsection{Limitações metodológicas}

Além das já mencionadas, destacam-se: (i) a amostra é pequena (15 observações nas regressões com PIB), o que reduz o poder dos testes e torna os resultados sensíveis a observações individuais; (ii) o PIB per capita varia muito entre os dois municípios e pouco ao longo do tempo dentro de cada um, de modo que seu coeficiente é estimado com pouca precisão quando se controla o efeito fixo de município; (iii) o PIB municipal é uma \textit{proxy} imperfeita da renda dos residentes, sobretudo em Manaus, onde o valor adicionado industrial da Zona Franca não se converte integralmente em renda local, e foi deflacionado pelo IPCA por falta de deflator municipal; (iv) a população de 2023 é interpolada; (v) a qualidade da DCA depende do envio pelos municípios, e não foi possível auditar os lançamentos individuais; e (vi) com dez anos de dados, os testes para a Espanha também têm baixo poder estatístico. Esses limites não são contornados com imputações ou com a inclusão de outros municípios: a análise se restringe às unidades atribuídas, e as conclusões são formuladas na medida do que esses dados permitem afirmar.

\section{Resultados e discussão}\label{sec:resultados}

\subsection{O que mudou em relação ao trabalho anterior}

O trabalho anterior\footnote{VICTOR, C. \textit{Prova P1 de Finanças do Setor Público} (matrícula 231020069). 2026. Repositório GitHub: \url{https://github.com/cauaviictor/P1---FSP---231020069---CAUA-VICTOR}, \textit{commit} 5209bcec.} disponibilizou base editável (planilha com fórmulas e arquivos do Finbra), preservada integralmente na pasta \texttt{01\_base\_original} do repositório. A conferência dessa base orientou as correções e ampliações resumidas no Quadro~\ref{tab:mudancas}; o registro detalhado está em \texttt{05\_registro\_alteracoes/registro\_alteracoes.md}.

\begin{table}[htbp]
\centering
\caption{Principais diferenças entre a base do trabalho anterior e a base atualizada}\label{tab:mudancas}
\scriptsize
\begin{tabular}{p{2.6cm}p{5.7cm}p{6.2cm}}
\toprule
Aspecto & Trabalho anterior (P1) & Esta versão (P2)\\
\midrule
Unidades & Espanha; Maceió/AL e Duque de Caxias/RJ & Espanha; Manaus/AM e Parintins/AM (redistribuição do professor)\\
Período & 2013--2022 & 2016--2025 (dez exercícios mais recentes com DCA completa)\\
Fonte municipal & Planilhas do Finbra baixadas manualmente (AL e RJ) & API do Siconfi (DCA, Anexos I-C, I-D e I-E), com coleta registrada\\
Plano de contas & Não tratado; contas lidas por código & Harmonização das três classificações (2016--17, 2018--21, 2022--25) pelo nome da conta, com regra anti-dupla contagem\\
Intraorçamentárias & Incluídas na despesa & Excluídas de receita e despesa\\
FUNDEB & Receita bruta, sem deduções & Receitas líquidas das deduções; FUNDEB recebido entre as transferências\\
Deflator & Fatores pontuais da Calculadora do Cidadão, base 2022, com fator de 1,0054 no próprio ano-base & Média anual do IPCA (Sidra 1737), base 2025 = 1; IPCH anual para a Espanha\\
População & Estimativas até 2021 e Censo 2022 na mesma série; números armazenados como texto & Série harmonizada aos Censos de 2010 e 2022; série oficial mantida para sensibilidade\\
País & FMI (WEO abr./2024), governo geral; ``equilíbrio estrutural'' rotulado de forma imprecisa & Eurostat por subsetor (central, regional, local, seguridade), saldo B.9 e dívida de Maastricht, até 2025\\
Dados ausentes & ``--'' misturado a números; IDEB ausente em anos pares & Ausências explícitas (NA), nunca zero; Parintins 2016 sem DCA registrado\\
Análises & Descritivas, Laffer, regressão simples receita-PIB & Regressões com efeito fixo de município para \textit{flypaper} e esforço tributário, com oito testes de sensibilidade e diagnósticos para amostra pequena; estudo bibliométrico\\
Reprodutibilidade & Planilha & Scripts Python e R que reproduzem bases, tabelas e figuras\\
\bottomrule
\end{tabular}
\par\smallskip
\begin{minipage}{0.97\textwidth}\scriptsize Fonte: elaboração própria a partir da conferência da base de Victor (2026).\end{minipage}
\end{table}

Algumas correções alteram a leitura dos números. A mais importante para a comparação temporal é a da população: misturar estimativas antigas com o Censo de 2022 faria a receita per capita de Parintins saltar artificialmente cerca de 20\% entre 2021 e 2022, apenas pela troca da base populacional. A exclusão das despesas intraorçamentárias reduz a despesa de Manaus em até R\$ 0,9 bilhão em um único ano (2020) e evita que contribuições patronais ao regime próprio sejam contadas como gasto final. Para a Espanha, a substituição do FMI pelo Eurostat corrige um problema conceitual: a variável chamada de ``equilíbrio estrutural'' na base anterior não corresponde ao saldo efetivo das administrações públicas, que em 2020 foi de \esBnoveVinte\% do PIB; a nova base usa a capacidade ou necessidade de financiamento (B.9), conceito empregado nas regras europeias.

\subsection{Manaus e Parintins: financiamento e gasto}

A Tabela~\ref{tab:focos} e a Figura~\ref{fig:composicao} apresentam a evolução das finanças dos dois municípios. Em Manaus, a receita total líquida real per capita passou de R\$ \recManIni{} em 2016 para R\$ \recManFim{} em 2025, crescimento de \recManVar\%. A receita tributária real per capita cresceu mais rapidamente (\tribManVar\%, de R\$ \tribManIni{} para R\$ \tribManFim), puxada pelo ISS, que responde em média por \issShareMan\% da arrecadação tributária, enquanto o IPTU responde por \iptuShareMan\%. Ainda assim, as transferências financiaram entre \depManMin\% e \depManMax\% da receita total no período (média de \depManMedia\%). A composição das transferências reflete a economia da capital: a cota-parte do ICMS, ligada ao valor adicionado da indústria incentivada, responde em média por \icmsShareMan\% das transferências líquidas, o FUNDEB por \fundebShareMan\% e o FPM por apenas \fpmShareMan\%.

\begin{figure}[htbp]
\centering
\caption{Receita real per capita por origem, Manaus e Parintins, 2016--2025 (R\$ de 2025)}\label{fig:composicao}
\includegraphics[width=0.95\textwidth]{fig1_composicao_receita_focos.png}
\par\smallskip
\begin{minipage}{0.95\textwidth}\scriptsize
Fonte: elaboração própria com dados do Siconfi/DCA e do IBGE. Receita corrente própria = receita corrente bruta menos transferências correntes brutas; transferências líquidas incluem as de capital e excluem as deduções para o FUNDEB. Parintins não tem DCA de 2016 no Siconfi.
\end{minipage}
\end{figure}

\begin{table}[htbp]
\centering
\caption{Manaus e Parintins: receitas, despesas e indicadores fiscais, 2016--2025}
\label{tab:focos}
\footnotesize
\begin{tabular}{lrrrrrrr}
\toprule
Ano & \makecell{Receita\\total líq.\\(R\$ pc)} & \makecell{Receita\\tributária\\(R\$ pc)} & \makecell{Transf.\\líquidas\\(R\$ pc)} & \makecell{Despesa\\empenhada\\(R\$ pc)} & \makecell{Resultado\\orçam.\\(R\$ mi)} & \makecell{Depen-\\dência\\(\%)} & \makecell{Rec. trib./\\RCL\\(\%)}\\
\midrule
\multicolumn{8}{l}{\textit{Manaus}}\\
2016 & 3.524,47 & 747,60 & 1.855,85 & 3.382,64 & 273,3 & 52,7 & 24,4\\
2017 & 3.180,48 & 741,13 & 1.999,21 & 3.252,93 & -141,2 & 62,9 & 23,5\\
2018 & 3.522,07 & 887,38 & 2.123,73 & 3.391,50 & 257,3 & 60,3 & 26,4\\
2019 & 3.944,53 & 929,00 & 2.223,45 & 3.826,82 & 234,6 & 56,4 & 26,9\\
2020 & 4.205,71 & 926,14 & 2.428,79 & 4.231,71 & -52,4 & 57,7 & 25,3\\
2021 & 4.178,39 & 954,71 & 2.574,70 & 3.806,60 & 757,9 & 61,6 & 25,5\\
2022 & 4.236,01 & 1.027,30 & 2.668,12 & 4.342,18 & -219,1 & 63,0 & 25,3\\
2023 & 4.609,38 & 1.089,72 & 2.582,93 & 4.493,47 & 250,9 & 56,0 & 26,1\\
2024 & 4.715,63 & 1.198,48 & 2.704,09 & 4.746,61 & -70,6 & 57,3 & 27,2\\
2025 & 4.757,67 & 1.262,24 & 2.739,52 & 4.647,34 & 254,2 & 57,6 & 28,0\\
\multicolumn{8}{l}{\textit{Parintins}}\\
2016 & \multicolumn{7}{c}{DCA não homologada no Siconfi (dado ausente, não imputado)}\\
2017 & 2.597,80 & 170,05 & 2.324,54 & 2.745,85 & -14,6 & 89,5 & 6,6\\
2018 & 3.100,85 & 154,08 & 2.759,83 & 3.128,78 & -2,7 & 89,0 & 5,3\\
2019 & 3.010,26 & 157,02 & 2.727,34 & 2.994,96 & 1,5 & 90,6 & 5,5\\
2020 & 3.789,42 & 162,95 & 3.503,11 & 3.788,54 & 0,1 & 92,4 & 4,6\\
2021 & 4.095,29 & 233,93 & 3.739,29 & 3.960,13 & 13,1 & 91,3 & 6,3\\
2022 & 4.264,52 & 259,50 & 3.806,43 & 4.460,43 & -18,9 & 89,3 & 6,5\\
2023 & 4.878,49 & 265,63 & 4.416,44 & 4.584,00 & 29,2 & 90,5 & 5,7\\
2024 & 6.058,96 & 306,47 & 5.207,27 & 6.227,16 & -17,1 & 85,9 & 5,6\\
2025 & 7.005,42 & 351,06 & 6.210,64 & 6.439,76 & 57,6 & 88,7 & 6,0\\
\bottomrule
\end{tabular}
\par\smallskip
\begin{minipage}{0.97\textwidth}\scriptsize
Fonte: elaboração própria com dados do Siconfi/DCA (STN) e do IBGE (Sidra). Valores em reais de 2025 (IPCA, média anual) e per capita pela população harmonizada (Seção~\ref{sec:tratamentos}). Receita exceto intraorçamentária, líquida das deduções para o FUNDEB; despesa empenhada exceto intraorçamentária. Dependência = transferências correntes e de capital líquidas / receita total líquida.
\end{minipage}
\end{table}


Parintins tem um padrão de outra natureza. As transferências responderam por \depParMin\% a \depParMax\% da receita total (média de \depParMedia\%), e a receita tributária correspondeu, em média, a apenas \autParMedia\% da receita corrente líquida, contra \autManMedia\% em Manaus. Mesmo esse percentual superestima a autonomia: a receita tributária inclui o imposto de renda retido na fonte dos servidores municipais, que pertence ao município por determinação constitucional, mas depende da própria folha de pagamento e não de uma base econômica local. O ISS representa \issSharePar\% da receita tributária e o IPTU, \iptuSharePar\%. As transferências de Parintins são mais diversificadas e mais vinculadas que as de Manaus: o FUNDEB responde por \fundebSharePar\%, o FPM por \fpmSharePar\%, os repasses do SUS por \susSharePar\% e o ICMS por \icmsSharePar\%. Entre 2017 e 2025, as transferências reais per capita cresceram \transfParVar\%, com forte aceleração em 2024 e 2025, impulsionada por repasses do SUS e por transferências de capital, e a despesa real per capita acompanhou esse movimento (\despParVar\%), passando de R\$ \despParIni{} para R\$ \despParFim.

A composição do gasto reforça a diferença. Em Manaus, a despesa real per capita com pessoal cresceu \pessoalManVar\% no período, menos do que a despesa total, enquanto saúde e educação cresceram \saudeManVar\% e \educManVar\%. Em Parintins, a despesa com pessoal cresceu \pessoalParVar\%, a saúde \saudeParVar\% e os investimentos, que partiam de valores muito baixos (R\$ \investParIni{} per capita em 2017), chegaram a R\$ \investParFim{} em 2025 --- uma expansão que coincide com o aumento das transferências de capital. Quanto ao resultado orçamentário, Manaus registrou superávit em \resPosMan{} dos \nAnosMan{} exercícios, com valores reais entre R\$ \resMinMan{} milhões (\anoResMinMan) e R\$ \resMaxMan{} milhões (\anoResMaxMan); Parintins registrou superávit em \resPosPar{} dos \nAnosPar{} exercícios disponíveis. Em ambos os casos, os resultados oscilam em torno do equilíbrio, como esperado de entes sujeitos à LRF e sem acesso relevante a crédito.

\subsection{Dependência de transferências: Manaus, Parintins e Espanha}

A Figura~\ref{fig:dependencia} compara a dependência de transferências dos dois municípios com a dos subsetores espanhóis. A distância entre Manaus e Parintins é estável ao longo de todo o período: a dependência de Parintins nunca ficou abaixo de \depParMin\%, enquanto a de Manaus oscilou entre \depManMin\% e \depManMax\%. A diferença não decorre de oscilações conjunturais, mas da estrutura econômica: o PIB per capita de Manaus em 2023 (R\$ \pibpcManFim) era cerca de três vezes e meia o de Parintins (R\$ \pibpcParFim), e é essa base que sustenta o ISS e a cota-parte do ICMS da capital. A primeira parte da hipótese H1 é, portanto, confirmada.

\begin{figure}[htbp]
\centering
\caption{Transferências recebidas em proporção da receita total: Manaus, Parintins e subsetores da Espanha, 2016--2025}\label{fig:dependencia}
\includegraphics[width=0.9\textwidth]{fig2_dependencia_transferencias.png}
\par\smallskip
\begin{minipage}{0.9\textwidth}\scriptsize
Fonte: elaboração própria com dados do Siconfi, do IBGE e do Eurostat. Para a Espanha, transferências recebidas = D.7 + D.9 sobre a receita total do subsetor (não consolidada). Parintins não tem DCA de 2016 no Siconfi.
\end{minipage}
\end{figure}

Na Espanha, as Comunidades Autônomas receberam em transferências, em média, \esDepCA\% de sua receita, e as entidades locais, \esDepLoc\% (Tabela~\ref{tab:espanha}). A dependência regional espanhola é, portanto, semelhante à de Manaus, mas muito inferior à de Parintins. Duas ressalvas qualificam a comparação. Primeiro, no SEC 2010 as parcelas cedidas do IRPF e do IVA são registradas como impostos regionais, embora as alíquotas e as bases sejam, em larga medida, definidas pelo Estado; uma classificação análoga à brasileira elevaria a dependência regional espanhola. Segundo, as entidades locais espanholas não financiam saúde e educação, que no Brasil absorvem a maior parte das transferências vinculadas. Ainda assim, a segunda parte de H1 é confirmada para Parintins e, parcialmente, para Manaus: os níveis subnacionais espanhóis têm autonomia tributária muito maior (impostos próprios equivalem a \esAutCA\% da receita regional e \esAutLoc\% da receita local, contra \autManMedia\% e \autParMedia\% da receita corrente líquida em Manaus e Parintins), e as entidades locais dependem menos de transferências do que ambos os municípios.

\begin{table}[htbp]
\centering
\caption{Espanha: indicadores fiscais por subsetor das administrações públicas, 2016--2025}
\label{tab:espanha}
\scriptsize
\begin{tabular}{lrrrrrrrr}
\toprule
 & \makecell{Receita\\\% PIB\\(média)} & \makecell{Despesa\\\% PIB\\(média)} & \makecell{Saldo\\\% PIB\\2016} & \makecell{Saldo\\\% PIB\\2025} & \makecell{Dívida\\\% PIB\\2016} & \makecell{Dívida\\\% PIB\\2025} & \makecell{Transf.\\recebidas/\\receita (\%)} & \makecell{Impostos\\próprios/\\receita (\%)}\\
\midrule
Governo geral & 40,7 & 45,0 & -4,2 & -2,4 & 102,0 & 100,5 & 4,9 & 57,0\\
Administração central & 18,4 & 21,5 & -2,5 & -2,0 & 89,9 & 92,5 & 11,0 & 78,5\\
Comunidades Autônomas & 16,0 & 16,5 & -0,8 & -0,4 & 24,7 & 20,2 & 58,9 & 35,4\\
Entidades locais & 6,5 & 6,1 & 0,6 & 0,3 & 2,9 & 1,2 & 44,3 & 47,4\\
Seguridade social & 14,7 & 15,7 & -1,5 & -0,3 & 1,5 & 8,1 & 15,8 & 0,0\\
\bottomrule
\end{tabular}
\par\smallskip
\begin{minipage}{0.97\textwidth}\scriptsize
Fonte: elaboração própria com dados do Eurostat (\texttt{gov\_10a\_main}, \texttt{gov\_10dd\_edpt1}, \texttt{nama\_10\_gdp}). Subsetores não consolidados entre si; transferências recebidas = D.7 + D.9 (inclui transferências intragovernamentais D.73 e D.92); impostos próprios = D.2 + D.5 + D.91, segundo o SEC 2010, que registra como receita regional a parcela cedida do IRPF e do IVA. Saldo = capacidade (+) ou necessidade ($-$) de financiamento (B.9).
\end{minipage}
\end{table}


O contraste mais nítido está na disciplina fiscal (Figura~\ref{fig:espanha}). O déficit das administrações públicas espanholas caiu de \esBnoveGerIni\% do PIB em 2016 para \esBnoveGerFim\% em 2025, depois de atingir \esBnoveVinte\% na pandemia, e a dívida bruta oscilou em torno de 100\% do PIB (\esGDGerIni\% em 2016, \esGDvinte\% em 2020 e \esGDGerFim\% em 2025), muito acima da referência de 60\%. O desequilíbrio, porém, concentra-se na administração central e na seguridade social. As Comunidades Autônomas mantiveram déficits pequenos (média de \esBnoveCAMedia\% do PIB) e reduziram sua dívida de \esGDCAIni\% para \esGDCAFim\% do PIB, e as entidades locais registraram superávit em \esSupLocAnos{} dos dez anos, com dívida que caiu de \esGDLocIni\% para \esGDLocFim\% do PIB. Esse resultado é coerente com H4 e com a evidência de que a autonomia tributária e as regras de estabilidade endurecem a restrição orçamentária subnacional \cite{arespa2022,asatryan2015}. Ele também se aproxima do que se observa em Manaus e Parintins, que, sob a LRF, oscilam em torno do equilíbrio. A diferença é que, na Espanha, a disciplina convive com autonomia tributária, ao passo que em Parintins ela convive com dependência quase total de transferências --- a combinação que \citeonline{rodden2002} associa a pressões por resgate quando as regras de endividamento se afrouxam.

\begin{figure}[htbp]
\centering
\caption{Espanha: saldo e dívida bruta por subsetor das administrações públicas, 2016--2025 (\% do PIB)}\label{fig:espanha}
\includegraphics[width=0.95\textwidth]{fig4_espanha_subsetores.png}
\par\smallskip
\begin{minipage}{0.95\textwidth}\scriptsize
Fonte: elaboração própria com dados do Eurostat (\texttt{gov\_10dd\_edpt1}). Dívida de Maastricht não consolidada entre subsetores.
\end{minipage}
\end{figure}

\subsection{Transferências, renda e gasto em Manaus e Parintins}

A Figura~\ref{fig:dispersao} mostra a relação entre transferências e despesa per capita nos dois municípios, ano a ano. Em Parintins, os pontos sobem praticamente sobre a reta de 45 graus: cada aumento das transferências se converteu em aumento quase equivalente da despesa. Em Manaus, os pontos ficam acima da reta e se espalham mais, porque a despesa também é financiada por receitas próprias, que cresceram no período. A comparação das variações acumuladas confirma essa leitura: entre \propDeMan{} e 2023, a despesa real per capita de Manaus aumentou R\$ \dDespMan{}, para um aumento de R\$ \dTransfMan{} nas transferências (R\$ \propTrMan{} de despesa por real transferido) e de R\$ \dPibMan{} no PIB per capita; em Parintins, entre \propDePar{} e 2023, a despesa aumentou R\$ \dDespPar{} para um aumento de R\$ \dTransfPar{} nas transferências (R\$ \propTrPar{} por real transferido) e de apenas R\$ \dPibPar{} no PIB per capita.

\begin{figure}[htbp]
\centering
\caption{Transferências líquidas e despesa empenhada per capita, Manaus e Parintins, 2016--2025 (mil R\$ de 2025)}\label{fig:dispersao}
\includegraphics[width=0.78\textwidth]{fig3_dispersao_transf_despesa.png}
\par\smallskip
\begin{minipage}{0.78\textwidth}\scriptsize
Fonte: elaboração própria com dados do Siconfi e do IBGE. Cada ponto é um ano (rótulo = dois últimos dígitos); as retas são ajustes lineares por município.
\end{minipage}
\end{figure}

\begin{table}[htbp]
\centering
\caption{Manaus e Parintins: transferências, renda local e gasto per capita, 2016--2023}
\label{tab:reg}
\scriptsize
\setlength{\tabcolsep}{2.6pt}
\begin{tabular}{lccccccccc}
\toprule
 & (1) & (2) & (3) & (4) & (5) & (6) & (7) & (8) & (9)\\
 & Níveis & Tendência & FPM & Log & Dif. & Pop. of. & Manaus & Parintins & Trib.\\
 & Desp. & Desp. & Desp. & ln Desp. & $\Delta$Desp. & Desp. & Desp. & Desp. & Trib.\\
\midrule
Transferências pc ($\beta_1$) & $0,899^{***}$ & $0,550$ &  & $0,839^{***}$ & $0,515$ & $1,040^{***}$ & $0,051$ & $1,054^{**}$ & $0,034$\\
 & (0,119) & (0,418) &  & (0,172) & (0,557) & (0,275) & (1,021) & (0,315) & (0,025)\\[1pt]
FPM pc &  &  & $1,064$ &  &  &  &  &  & \\
 &  &  & (1,040) &  &  &  &  &  & \\[1pt]
Demais transferências pc &  &  & $0,847^{***}$ &  &  &  &  &  & \\
 &  &  & (0,138) &  &  &  &  &  & \\[1pt]
PIB pc ($\beta_2$) & $0,051$ & $-0,024$ & $0,052$ & $0,042$ & $-0,038$ & $0,034$ & $0,129$ & $-0,108$ & $0,031^{***}$\\
 & (0,039) & (0,113) & (0,042) & (0,745) & (0,129) & (0,024) & (0,096) & (0,186) & (0,009)\\[1pt]
Parintins (efeito fixo) & $1,077$ & $-1,833$ & $1,125$ & $-0,300$ &  & $0,303$ &  &  & $0,556$\\
 & (1,661) & (4,463) & (1,846) & (0,955) &  & (1,052) &  &  & (0,386)\\[1pt]
Tendência linear &  & $0,145$ &  &  &  &  &  &  & \\
 &  & (0,178) &  &  &  &  &  &  & \\[1pt]
\midrule
$H_0$: $\beta_1=\beta_2$ ($t$) & 5,84$^{***}$ & 1,75 & 0,96 & 0,89 & 1,04 & 3,41$^{***}$ & -0,07 & 2,35$^{*}$ & --\\
Observações & 15 & 15 & 15 & 15 & 13 & 13 & 8 & 7 & 15\\
$R^2$ ajustado & 0,870 & 0,877 & 0,858 & 0,878 & 0,011 & 0,909 & 0,739 & 0,934 & 0,985\\
\bottomrule
\end{tabular}
\par\smallskip
\begin{minipage}{0.98\textwidth}\scriptsize
Fonte: elaboração própria (código \texttt{03\_analise.py}). MQO com dados empilhados de Manaus e Parintins (município-ano). Variáveis em milhares de reais de 2025 per capita, exceto (4), em logaritmos, e (5), em primeiras diferenças. A coluna (6) usa a população oficial do IBGE (sem 2023); as colunas (7) e (8) estimam cada município separadamente; a coluna (9) tem a receita tributária per capita como variável dependente. Na coluna (3), o teste compara o coeficiente do FPM com o do PIB. Erros-padrão robustos HC3 entre parênteses, com distribuição $t$ de Student. A amostra termina em 2023 porque o PIB municipal ainda não foi divulgado para 2024--2025, e Parintins não tem DCA de 2016. *** $p<0{,}01$; ** $p<0{,}05$; * $p<0{,}10$.
\end{minipage}
\end{table}


A Tabela~\ref{tab:reg} apresenta as estimações. Na especificação principal (coluna 1), com efeito fixo de município, um real adicional de transferências líquidas per capita associa-se a R\$ \bTrUm{} de despesa per capita, enquanto um real adicional de PIB per capita associa-se a R\$ \bPibUm{}, coeficiente não significativo. A hipótese de que as duas propensões são iguais é rejeitada ($t=\wtUm$, $p\,\wpUm$), resultado compatível com o efeito \textit{flypaper}. A ordem de grandeza está em linha com a literatura: a propensão a gastar a renda local fica na faixa de poucos centavos indicada por \citeonline{hines1995}, e a das transferências se aproxima dos valores encontrados para municípios brasileiros \cite{cossio2001,ferreira2019,ferreira2022} e em desenhos quase-experimentais em outros países \cite{dahlberg2008,leduc2017}. O resultado se mantém quando os valores per capita são calculados com a população oficial do IBGE (coluna 6: \bTrSeis{}, com $t=\wtSeis$ para a diferença) e quando se exclui o PIB para usar todos os anos até 2025 (\bTrOito{}, com \nOito{} observações).

A robustez, contudo, é limitada, como se espera com tão poucas observações. Com uma tendência linear comum (coluna 2), o coeficiente das transferências cai para \bTrDois{} e deixa de ser significativo, e a diferença em relação ao PIB passa a ter $p=\wpDois$; em primeiras diferenças (coluna 5), o coeficiente é de \bTrCinco{}, também sem significância. Isso indica que parte da associação em níveis decorre do crescimento simultâneo de transferências e despesas ao longo do tempo, e não apenas da resposta ano a ano. A estimação separada por município (colunas 7 e 8) mostra de onde vem o resultado: em Parintins, o coeficiente das transferências é de \bTrDez{} ($p=\pTrDez$), muito próximo de um real gasto por real recebido; em Manaus, é praticamente nulo (\bTrNove), porque a despesa da capital acompanha também a receita própria. Na especificação logarítmica (coluna 4), a elasticidade do gasto em relação às transferências é de \bTrQuatro{}, mas a do PIB é estimada com tanta imprecisão que a diferença entre as duas não é significativa. A coluna 3 separa o FPM das demais transferências: o coeficiente do FPM (\bFpmTres) é impreciso, porque o FPM per capita varia pouco ao longo do tempo dentro de cada município, enquanto o das demais transferências (\bOutTres) é significativo. A hipótese H2 é, portanto, apenas parcialmente confirmada: o padrão de \textit{flypaper} é claro em Parintins e na especificação em níveis, mas não resiste ao controle de tendência nem às primeiras diferenças.

A coluna 9 estima a Equação~\eqref{eq:trib}. O coeficiente das transferências sobre a receita tributária per capita é pequeno e não significativo (\bTrSete, $p=\pTrSete$), ao passo que o do PIB é positivo e significativo (\bPibSete, $p=\pPibSete$): a arrecadação própria acompanha a economia local, e não as transferências. Não há, portanto, evidência de que transferências maiores se associem a menor arrecadação própria, e a hipótese H3 não é confirmada. O resultado converge com a ausência de deslocamento da receita própria encontrada por \citeonline{litschig2013} e diverge da redução de eficiência arrecadatória documentada por \citeonline{mattos2011} e da substituição observada em municípios com \textit{royalties} \cite{postali2015}. Uma interpretação plausível é que, em um município como Parintins, com receita tributária próxima de \autParMedia\% da receita corrente líquida, não há muito esforço a ser substituído: a margem de ajuste está no gasto, e não na arrecadação.

Os diagnósticos do modelo da coluna 1 ajudam a qualificar esses resultados. Não há indício de heterocedasticidade (Breusch-Pagan com $p=\bpP$), de não normalidade dos resíduos (Shapiro-Wilk com $p=\swP$), de erro de forma funcional (RESET com $p=\resetP$) nem de autocorrelação relevante dos resíduos dentro de cada município (coeficiente AR(1) de \arRho, com $p=\arP$; Durbin-Watson de \dwMan{} em Manaus e \dwPar{} em Parintins). O VIF das transferências é baixo (\vifTr), mas os do PIB (\vifPib) e da variável de Parintins (\vifPar) são muito altos, porque o PIB per capita é quase inteiramente explicado pela diferença entre os dois municípios; é por isso que o coeficiente do PIB é impreciso. A observação mais influente é \cookMun{} em \cookAno{} (distância de Cook de \cookMax{}), ano de forte superávit da capital.

Para a Espanha, a regressão descritiva das despesas regionais e locais per capita sobre as transferências recebidas e o PIB per capita produz coeficientes de \esbTr{} e \esbPib{}, respectivamente ($n=\esN$), com diferença estatisticamente significativa. Com apenas vinte observações e transferências que, nas Comunidades Autônomas, financiam serviços de saúde e educação cujo gasto é em larga medida determinado pelo Estado, o resultado deve ser lido apenas como indício de que a mesma regularidade aparece em um sistema com regras fiscais mais estritas.

\subsection{Discussão: \textit{flypaper}, ilusão fiscal e os dois casos amazônicos}

Os resultados permitem responder à pergunta de pesquisa em três planos. No plano da dependência, Parintins financia cerca de nove de cada dez reais de sua receita com recursos de outras esferas, e mesmo Manaus, com a base própria mais ampla do estado, depende de transferências para mais da metade de sua receita. No plano do comportamento fiscal, o gasto de Parintins acompanha as transferências quase real a real e não responde à renda local, padrão compatível com o efeito \textit{flypaper}; em Manaus, a despesa responde também à receita própria, e o padrão é menos nítido. No plano comparativo, os níveis subnacionais espanhóis têm muito mais autonomia tributária e mantêm resultados próximos do equilíbrio, ao passo que Parintins mantém equilíbrio formal com autonomia mínima.

A ausência de associação entre transferências e esforço tributário, combinada com a forte associação entre transferências e gasto em Parintins, ajuda a discriminar entre as explicações da literatura. Ela é pouco compatível com a versão da ilusão fiscal em que os governos substituem tributos por transferências \cite{mattos2011}, mas é compatível com a versão em que a transferência reduz o preço percebido do bem público e expande o gasto \cite{logan1986,turnbull1998}, com o poder de agenda da burocracia \cite{filimon1982} e, sobretudo, com o argumento do custo da arrecadação: onde explorar a base local é caro, como em municípios insulares de baixa formalização, a resposta às transferências é maior \cite{aragon2013,dahlby2011}. As vinculações constitucionais (saúde, educação e FUNDEB) também aproximam mecanicamente o gasto das transferências \cite{cruz2020}, o que significa que parte da associação medida reflete regras de destinação, e não apenas escolhas locais.

Para a gestão pública, o caso de Manaus mostra que a dependência pode diminuir quando a base econômica sustenta tributos próprios --- a autonomia tributária subiu de \autManIni\% para \autManFim\% da receita corrente líquida no período ---, enquanto o caso de Parintins mostra o limite desse caminho em economias pequenas. A transição para o IBS, que substituirá o ISS e o ICMS, afetará justamente as duas fontes que distinguem Manaus de Parintins\footnote{BRASIL. Emenda Constitucional n.º 132, de 20 de dezembro de 2023 (Reforma Tributária). Disponível em: \url{https://www.planalto.gov.br/ccivil_03/constituicao/emendas/emc/emc132.htm}.}. A experiência espanhola indica que autonomia tributária e regras fiscais podem coexistir com descentralização ampla \cite{arespa2022,asatryan2015}, mas também que a transferência de impostos não elimina a necessidade de mecanismos de nivelamento \cite{zabalza2011}.

\subsection{Dados ausentes, sensibilidade e limitações dos resultados}

Os resultados devem ser lidos à luz de quatro limitações. A primeira é o tamanho da amostra: com 15 observações nas regressões com PIB, os testes têm baixo poder, e a conclusão sobre o efeito \textit{flypaper} depende da especificação, como mostram as colunas 2 e 5 da Tabela~\ref{tab:reg}. A segunda é a ausência da DCA de Parintins em 2016, que não foi imputada; as variações da série do município são calculadas a partir de 2017. A terceira é a cobertura do PIB municipal, que termina em 2023; a especificação sem PIB, que usa os anos até 2025, produz coeficiente semelhante ao da especificação principal. A quarta é a natureza associativa das estimativas: a endogeneidade das transferências e a presença de transferências vinculadas impedem interpretar $\beta_1$ como efeito causal \cite{knight2002}. Mudanças de conceito também afetam a leitura: as três classificações da receita foram harmonizadas pelo nome das contas, mas reclassificações pontuais (por exemplo, o tratamento do IRRF e das complementações da União ao FUNDEB) podem produzir pequenas descontinuidades, registradas no arquivo de auditoria.

\section{Conclusão}\label{sec:conclusao}

Este trabalho perguntou em que medida as transferências intergovernamentais condicionam o financiamento e o gasto de Manaus e Parintins entre 2016 e 2025 e se esse padrão é compatível com o efeito \textit{flypaper} e com a ilusão fiscal, à luz da organização fiscal subnacional da Espanha. A resposta é que as transferências condicionam fortemente ambos os municípios, em intensidades muito diferentes. Parintins financiou, em média, \depParMedia\% de sua receita com transferências, e sua receita tributária correspondeu a apenas \autParMedia\% da receita corrente líquida; Manaus, sustentada pela base econômica da Zona Franca, dependeu de transferências em \depManMedia\% da receita, com autonomia tributária crescente. Nas regressões com os dados anuais dos dois municípios, a despesa acompanha as transferências --- R\$ \bTrUm{} por real transferido --- e responde pouco à renda local --- R\$ \bPibUm{} por real de PIB per capita ---, diferença estatisticamente significativa na especificação principal e compatível com o efeito \textit{flypaper}. Esse padrão, porém, vem sobretudo de Parintins, onde a despesa cresce quase real a real com as transferências (\bTrDez{} por real), e perde significância quando se controla a tendência comum ou se usam primeiras diferenças; com tão poucas observações, ele deve ser lido como evidência sugestiva, e não conclusiva. Não se encontrou, contudo, associação entre transferências e arrecadação tributária, o que não sustenta a versão da ilusão fiscal baseada na substituição de esforço tributário.

Em relação às hipóteses, H1 (dependência maior em Parintins do que em Manaus e nos municípios do que nos subsetores espanhóis) foi confirmada para Parintins e, quanto à autonomia tributária, também para Manaus; H2 (propensão a gastar as transferências maior que a propensão a gastar a renda) foi confirmada parcialmente, de forma nítida em Parintins e na especificação em níveis, mas não nas especificações com tendência e em primeiras diferenças; H3 (redução do esforço tributário) não foi confirmada; e H4 (maior disciplina dos níveis subnacionais espanhóis do que da administração central) foi confirmada pelos dados do Eurostat, com Comunidades Autônomas próximas do equilíbrio e entidades locais superavitárias em nove de dez anos.

As contribuições do trabalho são três. A primeira é uma base de dados reprodutível, que preserva o trabalho anterior, documenta cada fonte e cada transformação e corrige problemas que afetavam diretamente as conclusões, como a mistura de bases populacionais, a inclusão de operações intraorçamentárias e o uso de um conceito inadequado de resultado para a Espanha. A segunda é empírica: o estudo traz para a literatura brasileira sobre efeito \textit{flypaper} dois municípios da Região Norte, pouco estudados, e mostra que, no município mais dependente, a associação entre transferências e gasto opera pela expansão da despesa, e não pela redução de tributos. A terceira é comparativa: o uso de dados por subsetor das contas nacionais espanholas permite comparar níveis de governo semelhantes, em vez de comparar municípios com o governo geral de um país.

As implicações para as finanças públicas são diretas. Em municípios com base tributária ínfima, como Parintins, a qualidade do gasto depende quase integralmente do desenho das transferências: regras de vinculação, critérios de rateio e mecanismos de prestação de contas são as principais alavancas de política. A transição para o IBS deve ser acompanhada de perto, porque redistribuirá justamente as receitas (ISS e cota-parte do ICMS) que diferenciam Manaus de Parintins. E a experiência espanhola sugere que ampliar a responsabilidade dos governos subnacionais pela arrecadação, combinada com regras fiscais críveis, é compatível com descentralização ampla.

As conclusões respeitam o alcance dos dados e do método. As estimativas medem associações condicionais com apenas 15 observações, e não efeitos causais; o PIB municipal é uma \textit{proxy} imperfeita da renda local e termina em 2023; a população de 2023 é interpolada; Parintins não tem DCA de 2016 no Siconfi; e a comparação com a Espanha é qualitativa, dadas as diferenças de classificação contábil e de competências. Pesquisas futuras podem ampliar a amostra para outros municípios e explorar as descontinuidades das faixas do FPM, com estimadores robustos à heterogeneidade dos efeitos \cite{callaway2021,dechaisemartin2020}, analisar separadamente transferências vinculadas e livres, incorporar indicadores de resultado das políticas (como o IDEB e indicadores de saúde) e acompanhar os efeitos da Reforma Tributária sobre a autonomia dos municípios amazônicos.


\bibliography{refs}

\appendix
\section*{Apêndice A -- Estudo bibliométrico: informações gerais, rede de coocorrência e mapa temático}
\addcontentsline{toc}{section}{Apêndice A -- Estudo bibliométrico}

\begin{table}[htbp]
\centering
\caption{Estudo bibliométrico: informações principais do corpus (\textit{bibliometrix})}
\label{tab:biblio}
\scriptsize
\begin{tabular}{@{}p{0.47\textwidth}@{\hspace{0.04\textwidth}}p{0.47\textwidth}@{}}
\begin{tabular}[t]{lr}
\toprule
\multicolumn{2}{l}{\textit{(a) Informações gerais}}\\
\midrule
Período coberto & 1952--2025\\
Documentos (artigos em periódicos) & 1992\\
Periódicos & 1109\\
Autores & 3107\\
Taxa de crescimento anual (\%) & 7,26\\
Citações médias por documento & 18,33\\
Coautores por documento & 1,94\\
Coautoria internacional (\%) & 8,333\\
\midrule
\multicolumn{2}{l}{\textit{(b) Periódicos mais produtivos}}\\
\midrule
Public Choice & 54\\
Public Finance Review & 28\\
International Tax and Public Finance & 26\\
Publius: the Journal of Federalism & 24\\
Public Budgeting \& Finance & 23\\
\bottomrule
\end{tabular}
&
\begin{tabular}[t]{lr}
\toprule
\multicolumn{2}{l}{\textit{(c) Países (afiliação dos autores, ocorrências)}}\\
\midrule
Indonésia & 551\\
Estados Unidos & 532\\
Brasil & 461\\
Ucrânia & 126\\
China & 110\\
Reino Unido & 85\\
\midrule
\multicolumn{2}{l}{\textit{(d) Palavras-chave mais frequentes}}\\
\midrule
intergovernmental transfers & 1094\\
fiscal federalism & 541\\
fiscal decentralization & 539\\
flypaper effect & 473\\
local government finance & 375\\
fiscal transfers & 368\\
fiscal illusion & 295\\
local government spending & 255\\
\bottomrule
\end{tabular}
\end{tabular}
\par\smallskip
\begin{minipage}{0.97\textwidth}\scriptsize
Fonte: elaboração própria com o pacote \textit{bibliometrix} \cite{aria2017} aplicado a registros do OpenAlex (consultas em título e resumo em inglês, português e espanhol, descritas na nota da Tabela~\ref{tab:baker}; apenas artigos com autoria e fonte primária classificada como periódico). Código: \texttt{05\_bibliometria.R}.
\end{minipage}
\end{table}


\begin{figure}[htbp]
\centering
\caption{Rede de coocorrência das 40 palavras-chave mais frequentes}\label{fig:rede}
\includegraphics[width=0.85\textwidth]{rede_coocorrencia.png}
\par\smallskip
\begin{minipage}{0.85\textwidth}\scriptsize
Fonte: elaboração própria com \textit{bibliometrix} (\texttt{biblioNetwork} e \texttt{networkPlot}, algoritmo de Fruchterman-Reingold), corpus do OpenAlex.
\end{minipage}
\end{figure}

\begin{figure}[htbp]
\centering
\caption{Mapa temático (centralidade $\times$ densidade) das palavras-chave}\label{fig:mapa}
\includegraphics[width=0.7\textwidth]{mapa_tematico.png}
\par\smallskip
\begin{minipage}{0.7\textwidth}\scriptsize
Fonte: elaboração própria com \textit{bibliometrix} (\texttt{thematicMap}, 250 termos, frequência mínima de 5).
\end{minipage}
\end{figure}

\section*{Apêndice B -- Repositório, reprodução e fontes}
\addcontentsline{toc}{section}{Apêndice B -- Repositório e reprodução}

O repositório que acompanha este artigo está organizado de acordo com os itens exigidos: \texttt{01\_base\_original} (base do trabalho anterior, sem alterações); \texttt{02\_dados\_brutos\_atualizados} (respostas das APIs do Siconfi, do IBGE, do Eurostat e do OpenAlex em JSON e as mesmas tabelas em CSV e XLSX editáveis); \texttt{03\_dados\_tratados} (bases usadas nas análises, em CSV e XLSX); \texttt{04\_dicionario\_variaveis} (conceito, unidade, fonte, endereço, período, cobertura, tratamento de ausentes e fórmula de cada variável); \texttt{05\_registro\_alteracoes} (comparação entre a base anterior e a nova versão); \texttt{06\_codigos\_e\_planilhas} (scripts Python e R e planilha com fórmulas que reproduzem tratamentos, cálculos, tabelas, gráficos e resultados); \texttt{07\_resultados} (tabelas, figuras e saídas do \textit{bibliometrix}); e \texttt{08\_artigo} (código-fonte deste artigo, em que todos os números do texto são importados automaticamente dos resultados). As instruções de reprodução estão no arquivo \texttt{README.md}.

Os endereços das fontes utilizadas são: Siconfi/DCA, \url{https://apidatalake.tesouro.gov.br/ords/siconfi/tt/dca}; IBGE/Sidra, \url{https://servicodados.ibge.gov.br/api/v3/agregados}; Eurostat, \url{https://ec.europa.eu/eurostat/api/dissemination/statistics/1.0/data/}; OpenAlex, \url{https://api.openalex.org}.


\end{document}
